\documentclass[11pt, a4paper]{article}

% Load custom definitions and packages
% Author: Dorian Thomé
% Description: Core Definitions and Packages for HEG Master Thesis

% ==========================================
% 1. Fonts & Language
% ==========================================
\usepackage[french,english]{babel}
\usepackage[utf8]{inputenc}
\usepackage[T1]{fontenc}
\usepackage{lmodern}  % Base font support
\usepackage{helvet}   % Loads Helvetica (Arial-like) 
\renewcommand{\familydefault}{\sfdefault} % Sets default text to Sans-Serif (Arial)
\usepackage{dsfont}
\usepackage{bbm}

% ==========================================
% 2. Spacing & Geometry
% ==========================================
\usepackage{setspace} 
\setstretch{1.2}      % 1.2 line spacing (Standard HEG requirement)
\setlength{\parindent}{0pt} % No indentation at the start of paragraphs
\setlength{\parskip}{6pt}   % A realistic gap between standard paragraphs 
\usepackage{microtype} % Better typography

\usepackage[a4paper, 
    left=3cm,         % Binding margin (3cm required)
    right=2cm,        % Standard margin (2cm)
    top=2cm,          
    bottom=2cm        
]{geometry}

% Create figures and schemas
\usepackage{tikz}
\usetikzlibrary{shapes.geometric, arrows.meta, positioning, calc, fit, backgrounds}
\usepackage{pgfplots}
\pgfplotsset{compat=1.18}

% ==========================================
% 3. Headers & Footers (Layout)
% ==========================================
\usepackage{fancyhdr}
\pagestyle{fancy}
\fancyhf{} % Clear all default headers/footers

% Remove the header line
\renewcommand{\headrulewidth}{0pt} 

\usepackage[hang]{footmisc}
\setlength{\footnotemargin}{1.25cm}

% ==========================================
% Custom Heading Sizes (matching Word template)
% ==========================================
\usepackage{titlesec}
\newcommand{\sectionbreak}{\clearpage}

% Level 1 Heading (16pt) - e.g., \section
% The {16}{19} means 16pt font size with 19pt line spacing.
\titleformat{\section}
  {\fontsize{16}{19}\bfseries\sffamily}{\thesection}{1em}{}

% Unnumbered Level 1 Heading (16pt) - Centered for Declaration, TOC, etc.
\titleformat{name=\section,numberless}
  {\fontsize{16}{19}\bfseries\sffamily\centering}{}{0em}{}

% Level 2 Heading (14pt) - e.g., \subsection
\titleformat{\subsection}
  {\fontsize{14}{17}\bfseries\sffamily}{\thesubsection}{1em}{}

% Level 3 Heading (12pt) - e.g., \subsubsection
\titleformat{\subsubsection}
  {\fontsize{12}{14}\bfseries\sffamily}{\thesubsubsection}{1em}{}

% Level 4 Heading (11pt) - e.g., \paragraph
\titleformat{\paragraph}
  {\fontsize{11}{13}\bfseries\sffamily}{\theparagraph}{1em}{}

% Note: If you are using the 'report' document class, you might also want to format \chapter to 16pt:
% \titleformat{\chapter}[hang]
%   {\fontsize{16}{19}\bfseries\sffamily}{\thechapter}{1em}{}

% Footer Configuration: 
% Left: 
%   Line 1: Report Title
%   Line 2: Author Name
% Right: 
%   Line 1: (empty)
%   Line 2: Page Number (Aligns with Author Name)
\lfoot{\fontsize{8}{10}\selectfont \shorttitle \\ \authorname} 
\rfoot{\fontsize{8}{10}\selectfont \\ \thepage}

% Increase footer space slightly to accommodate two lines
\setlength{\footskip}{30pt}

% Custom Double Light-Grey Footer Rule
\renewcommand{\footrule}{%
    \begingroup % Starts isolation
        \color{gray!50}% Set color to light grey
        \hrule width \headwidth height 0.5pt \vspace{0.5pt}% First line
        \hrule width \headwidth height 0.5pt \vspace{1.5pt}% Second line
    \endgroup % Ends isolation, protecting your main text
}

% ==========================================
% 4. Graphics, Tables & Math
% ==========================================
\usepackage{float}
\usepackage{array}
\usepackage{xcolor}
\usepackage{graphicx}
\usepackage{caption}
\usepackage{subcaption} 

% Force all figure and table captions to the top
\floatstyle{plaintop}
\restylefloat{figure}
\restylefloat{table}

\usepackage{tikz}
\usetikzlibrary{shapes.geometric, arrows.meta, positioning, calc, fit, backgrounds}
\usepackage{pgfplots}
\pgfplotsset{compat=1.18}

\usepackage{amsmath, amssymb, amsthm, amsfonts, amscd, amstext}
\usepackage{booktabs}
\usepackage{multirow}
\usepackage{siunitx}
\usepackage{enumitem}
\setlist{itemsep=0pt, parsep=0pt, topsep=6pt, partopsep=0pt}
\usepackage{csquotes}

% Rename Table of Contents
\addto\captionsenglish{%
  \renewcommand{\contentsname}{Table of contents}%
}

% ==========================================
% 5. Code Listings Style
% ==========================================
\usepackage{listings}
\lstdefinestyle{code}{
    basicstyle=\ttfamily\small,
    breaklines=true,
    backgroundcolor=\color{gray!10},
    frame=single,
    numbers=left,
    numberstyle=\tiny,
    keywordstyle=\color{blue},
    commentstyle=\color{gray!70},
    stringstyle=\color{orange!80!black}
}
\lstdefinestyle{llmprompt}{
    basicstyle=\ttfamily\small\color{black},
    backgroundcolor=\color{blue!5},
    frame=single,
    rulecolor=\color{blue!40},
    breaklines=true,
    breakatwhitespace=true,
    numbers=none
}

% ==========================================
% 6. References & Hyperlinks
% ==========================================
\usepackage[nottoc]{tocbibind}

\usepackage{hyperref}
\hypersetup{
    colorlinks=true,
    linkcolor=black,      % Black for print consistency
    citecolor=black,
    filecolor=black,
    urlcolor=black,
    pdfpagemode=FullScreen
}
\usepackage[english]{cleveref} % Must be loaded after hyperref

% --- Custom HEG Title Page ---
\newcommand{\makeHEGtitle}[5]{
% #1: Author Name
% #2: Supervisor Name(s)
% #3: Date/Place
% #4: HEG Logo Path
% #5: HES-SO Logo Path

\begin{titlepage}
    % 1. Logos (Absolute Position)
    \begin{tikzpicture}[remember picture, overlay]
        % Top Left: HEG Logo
        \node[anchor=north west, xshift=1cm, yshift=-1cm] at (current page.north west) {
            \includegraphics[width=0.3\textwidth]{#4}
        };
        % Bottom Right: HES-SO Logo
        \node[anchor=south east, xshift=-2.5cm, yshift=1.5cm] at (current page.south east) {
            \includegraphics[width=0.25\textwidth]{#5}
        };
    \end{tikzpicture}

    \begin{center}
        % Adjust this space to push the title down from the top logo
        \vspace*{3cm} 

        % 2. Title Section
        \begingroup 
            \linespread{1.5} \selectfont
            {\fontsize{18}{22}\selectfont \bfseries \fulltitle} \\[5cm] % Report Title Set to 18pt
        \endgroup

        % 3. Submission Statement & Author
        \begingroup 
            \linespread{1.2} \selectfont
            {\bfseries Master’s thesis by:}\\[0.5cm]
            {\bfseries #1} 
        \endgroup

        \vspace{2cm}

        % 4. Supervisor Section
        \begingroup 
            \linespread{1.2} \selectfont
            {Under the direction of:}\\[0.5cm]
            {\bfseries #2} % Supervisor Name
        \endgroup

        % 5. Push remaining text to bottom, but leave space for logo
        \vfill 

        % 6. Bottom Metadata
        \begingroup 
            \linespread{1.2} \selectfont
            % Date and Place
            {\bfseries #3}\\[2cm] 
            
            % School Information
            {\bfseries Information Sciences}\\
            {\bfseries Haute École de Gestion de Genève (HEG-GE)}
        \endgroup
        
        % Buffer to prevent text from hitting the bottom logo
        \vspace*{3.5cm} 
    \end{center}
\end{titlepage}
}

% Helper for the horizontal line
\newcommand{\hrue}[1]{\rule{\linewidth}{#1}}

\newenvironment{hegquote}
  {\list{}{\leftmargin=0.5cm \rightmargin=0.5cm \topsep=6pt \parsep=0pt}
   \item\relax\itshape\setstretch{1.0}}
  {\endlist}

% ==========================================
% Table of Contents Formatting (tocloft)
% ==========================================
\usepackage{tocloft}

% 1. Remove the extra vertical space before Level 1 titles (Sections)
\setlength{\cftbeforesecskip}{0pt}

% 2. Set Titles (Sections) to exactly 12pt (and keep them bold)
\renewcommand{\cftsecfont}{\fontsize{12pt}{14.4pt}\selectfont\bfseries}
\renewcommand{\cftsecpagefont}{\fontsize{12pt}{14.4pt}\selectfont\bfseries}

% 3. Set Subtitles (Subsections) to exactly 11pt AND make them bold
\renewcommand{\cftsubsecfont}{\fontsize{11pt}{13.2pt}\selectfont\bfseries}
\renewcommand{\cftsubsecpagefont}{\fontsize{11pt}{13.2pt}\selectfont\bfseries}

% 4. Set Sub-subtitles (Subsubsections) to exactly 11pt (normal weight)
\renewcommand{\cftsubsubsecfont}{\fontsize{11pt}{13.2pt}\selectfont}
\renewcommand{\cftsubsubsecpagefont}{\fontsize{11pt}{13.2pt}\selectfont}

% 5. Add dotted leaders to Level 1 titles (Sections)
\renewcommand{\cftsecdotsep}{\cftdotsep}

% Bibliography Configuration
\usepackage[backend=biber, style=iso-authoryear]{biblatex}
\addbibresource{../references.bib}
\renewcommand*{\bibfont}{\normalsize\setstretch{1.0}} % Forces 11pt and single spacing
\setlength{\bibitemsep}{12pt} % Sets the space between bibliography items

% --- Metadata ---
\newcommand{\authorname}{Dorian THOMÉ}
\newcommand{\shortauthors}{D. THOMÉ}
\newcommand{\fulltitle}{Bridging Reasoning and Reactive Execution in Multi-Agent Reinforcement Learning}
\newcommand{\shorttitle}{Bridging Reasoning \& Reactive Execution in MARL}
\newcommand{\place}{Geneva}
\newcommand{\keywords}[1]{\par \vspace{.01em}\noindent{\small\bfseries \keywordsname:} #1\par}
\def\keywordsname{Keywords}

\title{\fulltitle}
\author{\authorname}
\date{\today}

\begin{document}

% --- Generate Title Page ---
% Arguments: 
% 1. Author Name
% 2. Supervisor Name(s) (Use \\ for line breaks)
% 3. Place & Date
% 4. HEG Logo Path
% 5. HES-SO Logo Path
\makeHEGtitle
    {\authorname}
    {Alexandros KALOUSIS, Prof. Dr.}
    {Geneva, \today}
    {../images/heg_logo.png} 
    {../images/hesso_logo.png}

\pagenumbering{roman}
\section*{Declaration} % Required in the Word template
\addcontentsline{toc}{section}{Declaration}
This Master’s thesis is carried out as part of the final school examination, with a view to obtaining the title Master of Science HES-SO in Information Sciences.

The student attests that he has submitted the work to a plagiarism detection software. The student accepts the confidentiality clause, if applicable. The use of the conclusions and recommendations formulated in the Master’s thesis, without prejudice to their value, does not engage the responsibility of the author, the Master’s thesis, advisor, the jury or the HEG-GE.

\textit{“I hereby certify that this work has been carried out using only the sources cited in the bibliography, that it is the fruit of my personal reflection and has been written independently.”}

\vspace{1.5cm}
\hfill
\begin{tabular}{l}
Done at Geneva, 15th of August 2026 \\[0.5cm]
Dorian THOMÉ
\end{tabular}

\newpage

\section*{Acknowledgments}
\addcontentsline{toc}{section}{Acknowledgments}
I would like to express my deepest gratitude to my supervisor, Prof. Dr. Alexandros Kalousis, for his invaluable guidance and support throughout this research. I also thank M. Jacopo Castellini for his insightful feedback and guidance.

I acknowledge the use of the following resources:
\begin{itemize}
    \item Gemini 3.1 Pro for code generation, debugging, and assisting in spelling and grammar checks. Proof reading.
\end{itemize}

\newpage

\section*{Résumé}
\addcontentsline{toc}{section}{Résumé}

\begin{otherlanguage}{french}
Ce mémoire examine l'intégration des grands modèles de langage (LLM) avec l'apprentissage par renforcement multi-agents (MARL) afin de résoudre les goulets d'étranglement liés à l'exploration dans des environnements à récompenses rares. L'objectif principal est d'évaluer comment le raisonnement sémantique d'un LLM peut guider des agents optimisés par l'algorithme MAPPO à travers des tâches interdépendantes modélisées par un graphe orienté acyclique (DAG).

Pour ce faire, la méthodologie s'est articulée autour du développement et de l'évaluation de deux architectures neuro-symboliques distinctes. La première approche, le façonnage dynamique (Dynamic Reward Shaping), utilise le LLM comme un superviseur stratégique ajustant en continu les pondérations d'une fonction de récompense basée sur le potentiel (PBRS). La seconde approche introduit un méta-contrôleur hiérarchique (HRL) au sein duquel le LLM assigne des macro-actions discrètes (Options) aux agents, modifiant ainsi directement leur espace d'observation et leur signal de récompense intrinsèque. L'étude évalue l'impact de ces architectures à travers des études d'ablation ciblant le lissage temporel, l'adaptation sémantique en "zero-shot", l'ajustement fin (QLoRA) et la réflexion contrefactuelle.

Les résultats démontrent que l'architecture HRL surpasse largement le façonnage dynamique, permettant aux agents de maîtriser l'intégralité de l'environnement, y compris les phases de coordination complexes liées au combat. Le bilan critique des réalisations révèle que l'intégration de modèles sémantiques avec des politiques cinétiques nécessite des mécanismes de stabilisation stricts. L'utilisation d'une moyenne mobile exponentielle (EMA) et de garde-fous basés sur la divergence de Kullback-Leibler (KL) s'est avérée indispensable pour empêcher l'oubli catastrophique face à la non-stationnarité des objectifs. Par ailleurs, cette recherche met en lumière deux vulnérabilités fondamentales des systèmes neuro-symboliques. Premièrement, l'excès de confiance sémantique du LLM peut induire un piratage des récompenses (reward hacking), piégeant les agents dans des optima locaux. Deuxièmement, la réflexion contrefactuelle provoque une sur-correction induite face à de simples retards physiques, dégradant l'efficacité de l'apprentissage. D'un point de vue prospectif, bien que l'ajustement fin par QLoRA multiplie par cinq la vitesse de convergence, la viabilité de ces architectures pour le contrôle à haute fréquence nécessitera le développement de modèles de fondation dotés d'une véritable mémoire spatio-temporelle.

\keywords{Apprentissage par renforcement multi-agents (MARL), Grands modèles de langage (LLM), Apprentissage hiérarchique par renforcement (HRL), Façonnage de récompenses, Optimisation de la politique proximale (PPO), Systèmes neuro-symboliques, QLoRA.}
\end{otherlanguage}

\vspace{2cm}

\section*{Summary}
\addcontentsline{toc}{section}{Summary}
Reinforcement learning agents struggle in environments with sparse rewards. This thesis investigates combining Large Language Models (LLMs) with Multi-Agent Reinforcement Learning (MARL) to resolve exploration bottlenecks. The primary objective is to evaluate whether the semantic reasoning of an LLM can guide MAPPO agents through interdependent tasks modeled as a Directed Acyclic Graph (DAG).

We developed and evaluated two distinct neuro-symbolic architectures. The first approach, Dynamic Reward Shaping, utilizes the LLM as a strategic supervisor that continuously adjusts the priority weights of a Potential-Based Reward Shaping (PBRS) signal. The second approach introduces a Hierarchical Reinforcement Learning (HRL) Meta-Controller, where the LLM assigns discrete macro-actions (Options) to the agents, directly modifying their observation space and intrinsic rewards. Both architectures were evaluated through ablation studies targeting temporal smoothing, zero-shot semantic adaptation, QLoRA fine-tuning, and counterfactual reflection.

The results demonstrate that the HRL architecture significantly outperforms dynamic shaping, enabling agents to master the entire environment, including complex combat coordination. However, integrating semantic models with kinetic policies introduces severe non-stationarity. Implementing Exponential Moving Average (EMA) smoothing and Kullback-Leibler (KL) divergence guardrails proved necessary to prevent catastrophic forgetting. 

Furthermore, the research highlights two fundamental vulnerabilities in neuro-symbolic systems. First, the LLM's semantic overconfidence can induce reward hacking, trapping agents in local optima. Second, counterfactual reflection triggers over-correction to simple physical delays, actively degrading learning efficiency. Prospectively, while QLoRA fine-tuning accelerated convergence five-fold, the viability of these architectures for high-frequency control will require the development of foundation models equipped with spatio-temporal memory.

\textbf{Keywords:} Multi-Agent Reinforcement Learning (MARL), Large Language Models (LLMs), Hierarchical Reinforcement Learning (HRL), Reward Shaping, Proximal Policy Optimization (PPO), Neuro-Symbolic AI, QLoRA.

\newpage

% --- Frontmatter Lists ---
\begingroup
\setlength{\parskip}{0pt} % Temporarily remove paragraph spacing

\tableofcontents
\newpage

\listoftables 
\newpage

\listoffigures 
\newpage

\endgroup % Paragraph spacing returns to 6pt here

\pagenumbering{arabic}
\setcounter{page}{1} 

% Chapters
\section{Introduction}
\label{ch:introduction}

\subsection{Context: Multi-Agent Systems in Open-World Environments}
\label{sec:context}

Multi-Agent Reinforcement Learning (MARL) provides a mathematical framework for training decentralized agents to maximize a shared or competitive utility. As the field has matured, researchers have increasingly utilized Massively Multiplayer Online Role-Playing Games (MMORPGs) and open-world environments, such as Neural MMO and Minecraft, as testbeds for real-world complexity. These environments require agents to navigate vast state spaces, manage resources, and engage in high-level strategic coordination.

To solve these tasks, the Centralized Training with Decentralized Execution (CTDE) paradigm has emerged as the standard architecture. Under this paradigm, algorithms such as Multi-Agent Proximal Policy Optimization (MAPPO) \parencite{Yu_Velu_Vinitsky_Gao_Wang_Bayen_Wu_2022} have shown strong results in cooperative settings. By utilizing a centralized critic that has access to the global state during training, agents can learn robust decentralized execution policies that rely solely on local observations.

\subsection{The Exploration Bottleneck in Sparse-Reward Worlds}
\label{sec:problem_statement}

Despite the success of algorithms like MAPPO, pure deep reinforcement learning struggles catastrophically in environments characterized by long-horizon, sparse-reward structures. In complex RPG-style environments, progress is often gated behind strict dependency graphs (e.g., agents must jointly collect wood and stone, navigate to a workbench, craft a pickaxe, and then coordinate to mine iron). 

Traditional RL agents rely on random exploration to discover these initial reward signals. However, as the depth of the dependency chain increases, the probability of agents randomly executing the exact sequence of required motor commands approaches zero. While engineers traditionally attempt to solve this via dense, hand-crafted reward shaping, this introduces significant risks. Poorly designed shaping signals frequently lead to reward hacking, where agents optimize for surrogate metrics without achieving the true environmental milestone \parencite{ma2024eurekahumanlevelrewarddesign}. At the same time, as multiple agents update their policies simultaneously, the environment becomes highly non-stationary from the perspective of any individual agent, leading to variance spikes and policy collapse \parencite{Lowe_Wu_Tamar_Harb_Abbeel_Mordatch_2017}. The fundamental problem lies in bridging the gap between high-level semantic planning (knowing \emph{what} to do) and low-level kinetic control (learning \emph{how} to do it) without destabilizing the multi-agent learning process.

To systematically investigate this intersection of large language models and multi-agent reinforcement learning, this thesis poses the following three central research questions:

\begin{description}
    \item[\textbf{RQ1:}] How can a high-level, semantic language model be effectively integrated with high-frequency, low-level MARL agents to solve long-horizon tasks without destabilizing the training loop?
    \item[\textbf{RQ2:}] Between continuous dynamic reward shaping and discrete Hierarchical Reinforcement Learning (HRL), which architecture provides better sample efficiency and agent coordination?
    \item[\textbf{RQ3:}] What practical failure modes arise when coupling LLMs with MARL (such as semantic over-confidence or reactive self-debugging), and how can mathematical guardrails mitigate these risks?
\end{description}

\subsection{Proposed Approach: Neuro-Symbolic MARL}
\label{sec:proposed_approach}

To resolve this exploration bottleneck, this thesis proposes and evaluates two neuro-symbolic MARL architectures. We treat the learning process as a dual-tier problem, introducing Large Language Models (LLMs) as high-level semantic reasoners coupled with low-level MAPPO execution policies \parencite{Yu_Velu_Vinitsky_Gao_Wang_Bayen_Wu_2022}. This approach leverages the pre-trained world knowledge of LLMs to synergize reasoning and acting \parencite{Yao_Zhao_Yu_Du_Shafran_Narasimhan_Cao_2023}, bypassing the need for the agents to discover basic logical dependencies via trial-and-error.

Specifically, the subject is addressed through two implementations. The first, termed the Mixing Board, relies on dynamic reward shaping. The LLM acts as a meta-supervisor that observes aggregate training statistics and dynamically adjusts priority weights over subtasks. These weights govern a continuous Potential-Based Reward Shaping (PBRS) signal \parencite{Ng_Harada_Russell_1999}, which guides the MAPPO agents toward critical bottlenecks while theoretically preserving the optimal policy.

The second implementation, the Meta-Controller, transitions to a discrete hierarchical reinforcement learning architecture. Building upon the Options framework \parencite{Sutton_Precup_Singh_1999} and hierarchical meta-controllers \parencite{Kulkarni_Narasimhan_Saeedi_Tenenbaum_2016}, the LLM replaces the traditional symbolic planner by issuing temporally extended macro-actions directly to the agents. The agents' observation space is augmented with one-hot option encodings, and strict Kullback-Leibler (KL) divergence guardrails enforce trust region constraints during policy updates. This approach allows the agents to master the complete dependency graph and adapt to zero-shot semantic shifts in the environment. Through both architectures, this work demonstrates how the zero-shot planning capabilities of quantized, self-hosted language models (QLoRA) \parencite{Dettmers_Pagnoni_Holtzman_Zettlemoyer_2023} can be integrated into deep reinforcement learning pipelines to solve complex cooperative tasks.

\subsection{Constraints and Sustainability}
\label{sec:sustainability}

A critical constraint, and a core design philosophy, of this thesis is the reliance on localized, consumer-grade computing. The contemporary trend in neuro-symbolic AI heavily favors routing reasoning tasks to massive, proprietary models via commercial APIs. However, this paradigm introduces severe latency, raises privacy concerns, and incurs prohibitive financial costs when queried across millions of MARL training steps. Furthermore, it relies on foundation models with significant training energy costs \parencite{strubell2019energy}. 

To address these practical limitations, this work is explicitly constrained to operate entirely on localized hardware. By utilizing heavily quantized, self-hosted language models optimized via QLoRA, we demonstrate that complex, hierarchical multi-agent coordination can be achieved without cluster-scale computing. This commitment to ``Green AI'' \parencite{schwartz2020green} not only minimizes the energetic footprint of the training loop, but also proves that such architectures are viable for future deployment on resource-constrained edge devices, such as robotic swarms or autonomous systems operating on strict power budgets.

\subsection{Thesis Outline}
\label{sec:outline}

The subsequent chapters establish the theoretical foundations of Multi-Agent Reinforcement Learning and Large Language Models before detailing the experimental testbed. To rigorously evaluate the proposed architectures, the environment is structured as a strict task dependency graph, specifically, a Directed Acyclic Graph (DAG) where higher-level goals are gated by sequential milestones. Against this backdrop, the thesis presents its core contributions: the continuous dynamic reward shaping framework (Mixing Board) and the discrete Hierarchical Reinforcement Learning architecture (Meta-Controller). Following the architectural design, the experimental results are analyzed to compare sample efficiencies and isolate the critical failure modes of neuro-symbolic coupling. Finally, the work concludes with a discussion on future directions for sustainable, edge-viable MARL systems.

\section{Background: Multi-Agent Reinforcement Learning}
\label{ch:background_marl}

This chapter provides the necessary background on Multi-Agent Reinforcement Learning (MARL). 
Separating this from the LLM background helps keep the text focused and modular.

\subsection{Reinforcement Learning Fundamentals}

\subsection{Multi-Agent Systems}

\subsection{State of the Art in MARL}

\subsection{MAPPO Agents: Training and Gradient Flow}
In Multi-Agent Proximal Policy Optimization (MAPPO), the system relies on the Centralized Training with Decentralized Execution (CTDE) paradigm. This means there are two distinct neural networks (or sets of networks) being trained: the Actor and the Critic.

\subsubsection{The Actor Network (Policy)}
\textbf{What is trained:} The Actor network, parameterized by $\theta$, learns the policy $\pi_\theta(a_t\vert{}s_t)$. This defines the probability of an agent taking a specific action $a_t$ given its local observation $s_t$.

\textbf{The Objective Function:} The Actor is trained to maximize a clipped surrogate objective function. We calculate the probability ratio $r_t(\theta)$ between the new policy and the old policy:
\begin{equation}
r_t(\theta) = \frac{\pi_\theta(a_t\vert{}s_t)}{\pi_{\theta_{\text{old}}}(a_t\vert{}s_t)}
\end{equation}
The Actor's loss function mathematically prevents the policy from updating too aggressively in a single step (which would cause training collapse):
\begin{equation}
L^{\text{CLIP}}(\theta) = \hat{\mathbb{E}}_t \left[ \min\left( r_t(\theta) \hat{A}_t, \text{clip}(r_t(\theta), 1 - \epsilon, 1 + \epsilon) \hat{A}_t \right) \right]
\end{equation}
Here, $\hat{A}_t$ is the Advantage estimate, and $\epsilon$ is the clipping hyperparameter (usually 0.2).

\subsubsection{The Centralized Critic Network (Value)}
\textbf{What is trained:} The Critic network, parameterized by $\phi$, learns the value function $V_\phi(S_t)$. Crucially, $S_t$ is not just the local observation; it is the global joint state (or the concatenation of all agents' observations). This allows the Critic to accurately assess the true value of a state across the entire multi-agent system.

\textbf{The Objective Function:} The Critic minimizes the Mean Squared Error (MSE) between its predicted value and the actual calculated returns $\hat{R}_t$ (usually derived via Generalized Advantage Estimation, GAE):
\begin{equation}
L^{\text{VF}}(\phi) = \hat{\mathbb{E}}_t \left[ \left( V_\phi(S_t) - \hat{R}_t \right)^2 \right]
\end{equation}

\subsubsection{How the Gradients Flow in MAPPO}
During a training update:
\begin{itemize}
    \item \textbf{Critic Gradient Flow:} The error $L^{\text{VF}}(\phi)$ is computed. Gradients flow backward strictly through the Centralized Critic network, updating the weights $\phi$ to make its future predictions more accurate.
    \item \textbf{Actor Gradient Flow:} The Advantage $\hat{A}_t$ is computed using the Critic's evaluations. If $\hat{A}_t$ is positive (the action was better than expected), the Actor's loss $L^{\text{CLIP}}(\theta)$ pushes the gradient to increase the probability of that action. The gradient flows backward entirely through the local Actor network, updating weights $\theta$.
\end{itemize}
\section{Background: Large Language Models}
\label{ch:background_llm}

This chapter introduces Large Language Models (LLMs) and their capabilities, specifically in the context of reasoning and acting.

\subsection{Foundations of LLMs}

\subsection{LLMs as Agents}
(e.g., ReAct, Prompting strategies)

\subsection{State of the Art in LLM Integration with RL}

\subsection{LLM Fine-Tuning: QLoRA and Cross-Entropy}
To clarify an important technical detail regarding the language model side: the LLM does not use Binary Cross-Entropy. Because an LLM must predict the next token out of a massive vocabulary (often 32,000 to 128,000 possible tokens), it uses Categorical Cross-Entropy (specifically, Causal Language Modeling Loss).

\subsubsection{The Loss Function (Causal Language Modeling)}
\textbf{What is trained:} The LLM is trained to maximize the likelihood of the next token in a sequence, given all preceding tokens.

If we have a context sequence $X$ and a target sequence of tokens $y$, the loss function is the negative log-likelihood over the sequence of $N$ tokens:
\begin{equation}
L_{\text{LLM}} = -\frac{1}{N} \sum_{i=1}^{N} \log P(y_i \mid y_{<i}, X)
\end{equation}
The probability $P(y_i)$ is calculated by applying a Softmax function over the final output logits across the entire vocabulary dimension.

\subsubsection{Low-Rank Adaptation (LoRA) Formulation}
In modern LLM fine-tuning, we do not train the entire billion-parameter model (Full Fine-Tuning), as this requires massive GPU clusters. Instead, we use Parameter-Efficient Fine-Tuning (PEFT), specifically QLoRA (Quantized Low-Rank Adaptation).

\textbf{The Math of LoRA:} We freeze the pre-trained weight matrix $W_0 \in \mathbb{R}^{d \times k}$ and inject two tiny trainable matrices, $A \in \mathbb{R}^{r \times k}$ and $B \in \mathbb{R}^{d \times r}$, where the rank $r$ is very small (e.g., 8, 16, or 32). The forward pass becomes:
\begin{equation}
h = W_0 x + \Delta W x = W_0 x + \frac{\alpha}{r} B A x
\end{equation}
Here, $\alpha$ is a scaling constant.

\subsubsection{How the Gradients Flow in QLoRA}
This is where the computational efficiency happens:
\begin{itemize}
    \item \textbf{Forward Pass:} The input passes through the frozen 4-bit weights $W_0$ and the trainable 16-bit matrices $A$ and $B$.
    \item \textbf{Loss Calculation:} The output logits are compared against the target tokens using the Categorical Cross-Entropy loss $L_{\text{LLM}}$.
    \item \textbf{Gradient Flow (Backpropagation):} The gradients flow backward from the loss function, down through the network layers. However, when the gradient reaches the weight matrices, the gradients for $W_0$ are discarded (zeroed out). The optimizer only calculates and applies the gradients to the small matrices $A$ and $B$.
\end{itemize}

By freezing the base model and only routing gradients to the $A$ and $B$ adapters, memory consumption is reduced by up to 75\%, allowing you to fine-tune massive reasoning models on consumer hardware to act as your MARL meta-controller.

\section{Environment Design and Implementation}
\label{ch:the_environment}

Before detailing the neuro-symbolic architectures developed in this thesis, this chapter establishes the experimental testbed. Testing an LLM-guided MARL system requires an environment with both physical complexity (movement, spatial reasoning) and logical dependencies (tasks gated by preconditions). To this end, we engineered a custom 2D multi-agent grid-world environment from scratch.

\subsection{Environment Design \& Mechanics}
\label{sec:env_design}

The environment operates as a cooperative spatial grid where multiple agents must jointly gather resources, craft tools, and defeat an adversary to achieve a final objective. Unlike standard dense-reward environments, where agents receive incremental gradient signals based on Euclidean distance to a target, this environment utilizes a logical sparse-reward formulation. The agents receive extrinsic rewards solely upon the completion of discrete logical milestones (subtasks).

The overarching philosophy of this design is mathematical austerity. The environment acts as a consistent judge, providing no incremental guidance or continuous signals. It only issues a constant time penalty to encourage speed and milestone bonuses for completing predefined physical tasks.

\subsubsection{Extrinsic Reward Structure}
The environment reward $r_t^{\text{env}}$ is the true, unmodified signal from the environment. It consists of a constant time-step penalty of $-0.01$ to encourage speed, plus milestone bonuses upon successfully completing a physical subtask:

\begin{table}[H]
\centering
\caption{Environment reward structure from \texttt{crafting\_env.py}.}
\label{tab:env_rewards}
\begin{tabular}{@{}lc@{}}
\toprule
\textbf{Event} & \textbf{Reward (shared)} \\
\midrule
Step penalty      & $-0.01$ \\
Collect Wood      & $+2.0$  \\
Collect Stone     & $+2.0$  \\
Craft Pickaxe     & $+3.0$  \\
Mine Iron         & $+2.0$  \\
Craft Sword/Armor & $+3.0$  \\
Build Bridge      & $+3.0$  \\
Defeat Enemy      & $+10.0$ \\
Mine Gold (win)   & $+15.0$ \\
\bottomrule
\end{tabular}
\end{table}

\subsubsection{The Dependency Graph (DAG)}
\label{subsec:env_dag}

The fundamental exploration challenge of the environment is formulated as a Directed Acyclic Graph (DAG) of dependencies. Agents cannot arbitrarily navigate to the final goal. They must coordinate to satisfy a rigid sequence of material and infrastructural preconditions. The required topological task progression is illustrated in Figure \ref{fig:env_dag}.

\begin{figure}[H]
    \centering
    \begin{tikzpicture}[
        node distance=1.5cm,
        box/.style={rectangle, draw=black, thick, fill=blue!10, text width=3cm, align=center, rounded corners, minimum height=1cm},
        gather/.style={box, fill=green!10},
        craft/.style={box, fill=orange!10},
        combat/.style={box, fill=red!10},
        arrow/.style={-{Stealth[scale=1.2]}, thick, draw=gray!80}
    ]
        % Row 1
        \node[gather] (wood) at (3, 0) {1a. Collect Wood (x2)};
        \node[gather] (stone) at (7, 0) {1b. Collect Stone};

        % Row 2
        \node[craft] (bridge) at (0, -2) {2a. Build Bridge};
        \node[craft] (pickaxe) at (5, -2) {2b. Craft Pickaxe};
        
        % Row 3
        \node[gather] (iron) at (5, -4) {3. Mine Iron (x2)};
        
        % Row 4
        \node[craft] (sword) at (3, -6) {4a. Craft Sword};
        \node[craft] (armor) at (7, -6) {4b. Craft Armor};

        % Row 5
        \node[combat] (enemy) at (3, -8.5) {5. Defeat Enemy};
        
        % Row 6
        \node[gather] (gold) at (3, -10.5) {6. Mine Gold};

        % Edges Row 1 -> 2
        \draw[arrow] (wood.west) -| (bridge.north);
        
        \coordinate (pickaxe_join) at (5, -1.2);
        \draw[thick, draw=gray!80, rounded corners=4pt] (wood.east) -| (pickaxe_join);
        \draw[thick, draw=gray!80, rounded corners=4pt] (stone.west) -| (pickaxe_join);
        \draw[arrow] (pickaxe_join) -- (pickaxe.north);

        % Edges Row 2 -> 3
        \draw[arrow] (pickaxe.south) -- (iron.north);
        
        % Edges Row 3 -> 4
        \coordinate (iron_split) at (5, -4.8);
        \draw[thick, draw=gray!80, rounded corners=4pt] (iron.south) -- (iron_split);
        \draw[arrow] (iron_split) -| (sword.north);
        \draw[arrow] (iron_split) -| (armor.north);
        
        % Edges to Enemy
        \draw[arrow] (bridge.south) |- (enemy.west);
        \draw[arrow] (sword.south) -- (enemy.north);
        \draw[arrow] (armor.south) |- (enemy.east);
        
        % Edge to Gold
        \draw[arrow] (enemy.south) -- (gold.north);

    \end{tikzpicture}
    \caption[Environment Dependency DAG]{The strict dependency Directed Acyclic Graph (DAG) governing environment progression. Agents must jointly acquire basic resources to clear logical bottlenecks. Notably, Wood is required for both the Pickaxe and the Bridge, while Iron supplies both Sword and Armor.}
    \label{fig:env_dag}
\end{figure}

As the topological depth of this chain increases, the probability of agents discovering the complete sequence through random exploration decays. This structure isolates the exploration bottleneck that our neuro-symbolic models are designed to address.

\subsubsection{Action and Observation Spaces}
\label{subsec:env_spaces}

\subsection{Partial Observability and the Action Space}
To simulate realistic constraints, the agents do not possess global vision of the grid. Instead, they operate under strict partial observability. At each time step $t$, agent $i$ receives a localized observation vector $o_t^{(i)} \in \mathbb{R}^{15}$ (or $\mathbb{R}^{25}$ in Chapter 6), which includes: their own $(x,y)$ coordinates and inventory status. The environment hides target locations (e.g., resources, workbenches) from the state vector. This requires agents to explore the spatial grid to locate objectives. To further prevent the policies from overfitting to a static sequence of movements, the agents' starting coordinates are randomized at the beginning of every episode.

The discrete action space consists of five parameters: four cardinal movement directions (Up, Down, Left, Right) and one semantic ``Interact'' action. The interaction mechanic is context-sensitive. For instance, triggering the action while adjacent to a forest tile chops wood, whereas executing it at a workbench (contingent upon possessing the requisite inventory) initiates tool crafting.

\subsubsection{Visual Representation}
\label{subsec:env_visual}

Figure~\ref{fig:env_screenshots} provides a visual representation of the discrete grid-world environment utilized throughout the MARL training process. The interface renders the underlying state matrix, allowing human observers to track the spatial coordinates of the autonomous agents. While the core simulation runs headlessly to maximize throughput, this renderer serves as a tool for visual debugging and policy interpretation. 

\begin{figure}[H]
    \centering
    \includegraphics[width=\textwidth]{../images/capture_from_the_game.png}
    \vspace{1ex}
    {\raggedleft \footnotesize Screenshot by \authorname~2026. Sprites generated by Gemini 1.5 Pro, Google. Prompt used: ``generate each sprites with different state if needed like enemy alive and dead and make sure we can distinguish each elements'' [generated July 10, 2026]. \par}
    \caption[Visual Rendering of the MARL Environment]{Visual rendering of the custom MARL environment. The interface explicitly renders the underlying state matrix, allowing human observers to track the spatial coordinates of the autonomous agents as they navigate the map.}
    \label{fig:env_screenshots}
\end{figure}

All visual assets utilized in this rendering engine were generated using the Gemini 1.5 Pro model via the Antigravity coding assistant platform.

\subsection{Architecture \& Technology Stack}
\label{sec:env_tech}

Engineering a custom MARL environment demands a careful selection of the underlying technology stack, requiring a balance between visual debugging capabilities and simulation throughput.

Why Not a Game Engine? Initially, open-source game engines such as Godot~\parencite{godot} were considered for the simulation backbone. Godot provides a visual editor, collision physics, and a node-based hierarchy, making it an option for traditional game development.

The final constraint in our environment design is the hardware profile. Because MARL requires processing thousands of physical state transitions per second to calculate gradients, introducing an LLM query directly into the training loop introduces significant computational overhead.

As a result, the environment was engineered natively in Python, utilizing NumPy~\parencite{numpy} as the core mathematical engine and Pygame~\parencite{pygame} as a lightweight visual wrapper. This matrix-first approach offers architectural advantages for reinforcement learning:
\begin{itemize}
    \item \textbf{Direct Tensor Integration:} The underlying grid, entities, and inventory states exist purely as NumPy matrices. This avoids the serialization and deserialization overhead typically encountered when passing state data between a Python-based RL training script and an external C++ game engine.
    \item \textbf{Seamless NumPy-to-Game Translation:} Because the environment state is natively represented as a matrix, Pygame can translate these numerical coordinates into a visual interface for human debugging. The mathematical logic (kinematics, collisions) is decoupled from the rendering loop.
    \item \textbf{Headless Vectorization:} During the actual training loop, the Pygame visual renderer is disabled. The PyTorch models interact directly with the raw NumPy arrays, allowing a vectorized wrapper to step thousands of environments simultaneously in parallel.
\end{itemize}

\subsection{Libraries \& Infrastructure}
\label{sec:env_libs}

The codebase relies on a modern, standardized stack of Python libraries to seamlessly interface the custom environment with the deep neural networks and the large language models:

\begin{itemize}
    \item \textbf{MARL Interface:} The environment conforms to the PettingZoo API \parencite{Terry_Black_Grammel_Jayakumar_Hari_Sullivan_Santos_Perez_Horsch_Dieffendahl_et_al._2021}, which has established itself as the standard for multi-agent RL (analogous to the Gymnasium API for single-agent RL). Utilizing the Parallel API variant ensures the environment integrates natively with synchronous MAPPO tensor rollouts.
    \item \textbf{Deep Learning Framework:} PyTorch \parencite{Paszke_Gross_Massa_Lerer_Bradbury_Chanan_Killeen_Lin_Gimelshein_Antiga_et_al_2019} serves as the computational backbone for designing, compiling, and optimizing the MAPPO Actor and Critic networks, leveraging GPU acceleration for Generalized Advantage Estimation (GAE) and gradient descent.
    \item \textbf{LLM Integration:} To interface with the semantic Meta-Controllers, the Transformers library by Hugging Face is deployed \parencite{Wolf_Debut_Sanh_Chaumond_Delangue_Moi_Cistac_Rault_Louf_Funtowicz_et_al._2020}. For the Hierarchical Reinforcement Learning phase (which necessitates rapid, localized inference), the language models are quantized using bitsandbytes (QLoRA) \parencite{Dettmers_Pagnoni_Holtzman_Zettlemoyer_2023} and served via the Ollama framework~\parencite{ollama_2024} to strictly bound the latency introduced during the training loop.
    \item \textbf{Experiment Tracking:} To monitor the highly stochastic training process, TensorBoard is integrated directly into the PyTorch training loop. This allows for real-time visualization of critical metrics such as the episodic return, actor-critic loss curves, and policy entropy across millions of environment transitions.
\end{itemize}

\subsection{Hardware and Inference Architecture}
\label{sec:env_hardware}

A critical bottleneck in neuro-symbolic MARL is the severe latency introduced by the LLM during the training loop. This thesis relies on two distinct inference paradigms, online cloud APIs and localized serving, dictated by the constraints of the specific algorithms.

For the Dynamic Reward Shaping architecture, which requires infrequent, highly complex offline generation, this study leveraged the Gemini 2.5 Flash model hosted via Google AI Studio\footnote{https://aistudio.google.com/}. However, integrating cloud-based models directly into the high-frequency Hierarchical Reinforcement Learning (HRL) training loop proved intractable. The MAPPO environment requires thousands of synchronous environment steps per episode. Utilizing an online model for the HRL orchestrator generated an excessive volume of requests, rapidly exhausting API rate limits and introducing network latency that stalled the GPU-bound RL rollouts.

To overcome this, the HRL training loop was executed entirely locally on a consumer-grade workstation equipped with an NVIDIA RTX 5070 GPU (12GB VRAM). Running a foundation model locally alongside the PyTorch MAPPO actor-critic tensors within a strict 12GB memory envelope necessitated specific architectural compromises. To address this, the Ollama framework~\parencite{ollama_2024} was deployed to provide highly optimized, low-latency local inference.

As detailed in Section~\ref{sec:qwen_architecture}, the Qwen 2.5 7B architecture was selected specifically to satisfy this strict 12GB memory envelope. When combined with 4-bit quantization via Ollama, it fits comfortably alongside the MARL experience replay buffers without displacing the PyTorch actor-critic tensors.
\section{LLM Dynamic Implementation}
\label{ch:dynamic_implementation}

This chapter describes how the LLM is dynamically implemented and utilized within the system, detailing the reward generation or strategic guidance.

\subsection{System Architecture}

\subsection{Dynamic Prompting}

\subsection{Implementation Details}

% ==========================================================================
% Section 6 — HRL Architecture: The "Meta-Controller" Approach
% ==========================================================================
\section{Architecture 2: Hierarchical RL with an LLM Meta-Controller}
\label{ch:hrl_architecture}

The previous chapter established the \emph{Mixing Board} architecture, representing a "soft integration" of Large Language Models into MARL. In that paradigm, the LLM acted as an environmental architect, manipulating the gravitational pull of potential-based rewards to gently guide the agents toward logical milestones. While this separation preserves the underlying Markov properties, it restricts the LLM to an indirect role: it can only suggest a direction, not enforce a strategy. 

To fully explore the spectrum of neuro-symbolic integration, we must ask: what happens when we grant the LLM direct semantic authority over the agents' cognitive state?

This section details the second integration strategy: a Hierarchical Reinforcement Learning (HRL) framework. Here, we transition from soft environmental guidance to hard architectural integration. The LLM ceases to be a distant architect and becomes a direct \emph{Meta-Controller}. Instead of adjusting continuous weights over a fixed greedy planner, the Meta-Controller replaces the planner entirely. It selects discrete high-level \emph{Options} (macro-actions) that strictly govern the behavior of the low-level MAPPO policy.

In our established metaphor, the Commander is no longer just shifting the gravity; it is issuing explicit, discrete commands directly to the sensory inputs of the lumberjack and miner.

The exposition follows the implementation order:
\begin{enumerate}
    \item \textbf{Option Definitions} (\cref{sec:option_defs}): Defining the discrete macro-actions available to the Commander.
    \item \textbf{Discrete Option Selection} (\cref{sec:discrete_selection}): How the LLM is triggered to assign these explicit tasks.
    \item \textbf{One-Hot Observation Augmentation} (\cref{sec:onehot_aug}): How the Commander's text-based orders are mathematically injected into the MAPPO neural network.
    \item \textbf{Bipartite Reward Structure} (\cref{sec:bipartite_reward}): The dual-reward system required to train the workers to follow these new explicit orders.
    \item \textbf{Staleness and Counterfactual Reflection} (\cref{sec:lifecycle_guardrails}): How the Commander diagnoses its own mistakes when an order fails on the ground.
\end{enumerate}

% ──────────────────────────────────────────────────────────────────────────
\subsection{System Overview}
\label{sec:hrl_overview}

\Cref{fig:hrl_system_flow} illustrates this two-level command architecture. The Meta-Controller (LLM) observes inventory transitions and assigns explicit Options to each agent asynchronously. Unlike the previous approach where the MAPPO policy remained blind to the LLM's existence, here, the low-level MAPPO policy receives the active option directly as part of its mathematical observation space. This creates a true, tightly coupled two-level decision hierarchy.

\begin{figure}[H]
\centering
\begin{tikzpicture}[
    box/.style={draw, rounded corners, minimum width=2.6cm, minimum height=0.9cm,
                font=\small\sffamily, align=center, fill=gray!8},
    llmbox/.style={draw, rounded corners, minimum width=2.6cm, minimum height=0.9cm,
                font=\small\sffamily, align=center, fill=blue!8},
    arr/.style={-{Stealth[length=2.5mm]}, thick},
    node distance=1.6cm and 2cm
]
    % Top row: LLM level
    \node[llmbox] (llm) {LLM\\Meta-Controller};
    \node[box, right=3cm of llm] (ctrl) {OptionController\\$o_t^{(0)}, o_t^{(1)}$};

    % Bottom row: MAPPO level
    \node[box, below=2cm of llm] (env) {Env Step\\$s_t \to s_{t+1}$};
    \node[box, right=3cm of env] (agent) {MAPPO Agent\\$\pi_\theta(a \mid s, o)$};

    % LLM -> Controller
    \draw[arr] (llm) -- node[above,font=\footnotesize]{JSON} (ctrl);

    % Controller -> Agent (one-hot)
    \draw[arr] (ctrl.south) -- ++(0,-0.6) -|
        node[near start,right,font=\footnotesize]{one-hot $\mathbf{e}_{o}$} (agent.north);

    % Agent -> Env
    \draw[arr] (agent.south) -- ++(0,-0.4) -|
        node[near start,below,font=\footnotesize]{$a_t$} (env.south);

    % Env -> Agent (reward)
    \draw[arr] (env.east) -- node[below,font=\footnotesize]{$r_t^{\text{env}} + r_t^{\text{intr}}$} (agent.west);

    % Env -> LLM (trigger)
    \draw[arr,dashed] (env.north) -- node[left,font=\footnotesize]{\scriptsize success / stale} (llm.south);

    % Env -> Controller (reset on terminal)
    \draw[arr,dashed] (env.north east) -- node[right,font=\footnotesize,pos=0.3]{\scriptsize terminal} (ctrl.south west);
\end{tikzpicture}
\caption[Two-Level HRL Architecture]{Two-level HRL architecture. The LLM Meta-Controller is triggered by option success or staleness (dashed). The OptionController passes one-hot embeddings to the MAPPO agent at every step (solid), directly altering the agent's cognitive state.}
\label{fig:hrl_system_flow}
\end{figure}


% ──────────────────────────────────────────────────────────────────────────
\subsection{Option Definitions}
\label{sec:option_defs}

In the dynamic shaping approach, the LLM merely influenced a hard-coded planner. In this Hierarchical Reinforcement Learning (HRL) architecture, the LLM completely replaces the planner. It selects from a finite set of explicit macro-actions, or \emph{Options}, and directly commands the low-level MAPPO agents.

An Option $o\in\mathcal{O}$ is mathematically defined by the triple $\langle \mathcal{I}_o, \pi_o, \beta_o \rangle$: an initiation set, an intra-option policy, and a termination condition~\parencite{Sutton_Precup_Singh_1999}.
In our framework, the option set $\mathcal{O}$ is derived directly from the dependency DAG, providing 10 discrete macro-actions as detailed in Table~\ref{tab:option_props}.

\begin{table}[H]
\centering
\caption{Option properties: target zone, responsible agent, and initiation condition.}
\label{tab:option_props}
\small
\begin{tabular}{@{}lccc@{}}
\toprule
\textbf{Option} & \textbf{Zone} & \textbf{Agent} & \textbf{Initiation Condition} \\
\midrule
\texttt{COLLECT\_WOOD}  & Wood (0)      & 0     & Always \\
\texttt{COLLECT\_STONE} & Stone (1)     & 1     & Always \\
\texttt{CRAFT\_PICKAXE} & Workbench (2) & Both  & wood $\ge 1$, stone $\ge 1$, pickaxe $< 1$ \\
\texttt{MINE\_IRON}     & Iron (3)      & 1     & pickaxe $\ge 1$ \\
\texttt{CRAFT\_SWORD}   & Workbench (2) & Both  & iron $\ge 1$, sword $< 1$ \\
\texttt{CRAFT\_ARMOR}   & Workbench (2) & Both  & iron $\ge 1$, armor $< 1$ \\
\texttt{BUILD\_BRIDGE}  & Bridge (4)    & Both  & wood $\ge 1$, bridge $< 1$ \\
\texttt{FIGHT\_ENEMY}   & Enemy (5)     & Both  & bridge, sword, armor $\ge 1$ \\
\texttt{COLLECT\_GOLD}  & Gold (6)      & 1     & enemy defeated, gold $< 1$ \\
\texttt{IDLE}           & ---           & Both  & Always \\
\bottomrule
\end{tabular}
\end{table}

Rather than building 10 separate neural networks, the intra-option policy $\pi_o$ is executed by the single, shared MAPPO network, which conditions its kinetic movements based on the active option (detailed in Section~\ref{sec:onehot_aug}).

The termination condition ($\beta_o$) is not based on spatial coordinates, but on semantic success. For example, if the Commander radios the \texttt{COLLECT\_WOOD} Option to the lumberjack, that Option only terminates when the system registers an increase in the wood inventory.

% ──────────────────────────────────────────────────────────────────────────
\subsection{Discrete Option Selection by the LLM}
\label{sec:discrete_selection}

Because the Meta-Controller issues explicit, discrete commands, it does not need to continuously monitor the mathematical variance of the Critic network (unlike the Mixing Board approach). Instead, the Commander acts reactively, stepping in only when a worker completes a task or fails to make progress.

\subsubsection{Trigger Conditions: Success and Staleness}

Option selection is triggered dynamically under two specific conditions, evaluated per-en\-vi\-ron\-ment:

\begin{enumerate}
    \item \textbf{Success:} The agent successfully acquired the item (the inventory slot changed), meaning the current Option has terminated and the agent awaits a new order.
    \item \textbf{Staleness:} The active option has been running for more than $T_{\text{stale}}=150$ steps without completion.
    In MARL, a step $t \to t{+}1$ represents a physical movement or action in the grid. If an agent executes 150 low-level kinetic actions without satisfying its macro-objective, the system assumes the agent is stuck (e.g., trying to mine iron without a pickaxe) and forces an LLM intervention.
\end{enumerate}

Environments that are already awaiting an LLM response (\texttt{llm\_pending}) or are on a brief 50-step cooldown are safely excluded from triggering new queries.

\subsubsection{Prompt Structure}

When a trigger fires, the system translates the grid-world's physical inventory state into a text-based JSON schema.
For a standard \emph{Success} trigger, the Commander receives a Chain-of-Thought prompt~\parencite{Wei_Wang_Schuurmans_Bosma_Ichter_Xia_Chi_Le_Zhou_2022} detailing the current inventory and the DAG dependencies, requesting exactly one discrete Option per agent:

\begin{hegquote}
\begin{verbatim}
You are the high-level Cognitive Orchestrator for two MARL agents.
Strict Dependency DAG: (Wood+Stone) -> Pickaxe -> Iron
    -> (Sword+Armor) -> Bridge -> Enemy -> Gold.

### CURRENT STATE:
Inventory: {inventory}
Agent 0 Status: {a0_status}
Agent 1 Status: {a1_status}

### YOUR TASK:
Assign exactly ONE discrete Option to each agent.
Available Options: ["COLLECT_WOOD", "COLLECT_STONE", "CRAFT_PICKAXE",
    "MINE_IRON", "CRAFT_SWORD", "CRAFT_ARMOR", "BUILD_BRIDGE",
    "FIGHT_ENEMY", "COLLECT_GOLD", "IDLE"]

Respond ONLY with valid JSON:
{
  "dag_check": "<1 sentence reasoning>",
  "agent_0_option": "<Option>",
  "agent_1_option": "<Option>"
}
\end{verbatim}
\end{hegquote}

\subsubsection{Asynchronous Query Management}

Training MARL across 128 parallel environments generates significant query volume. To optimize throughput, the system deduplicates queries by grouping environments with identical inventory states, sending a single prompt per unique state to the LLM, and broadcasting the response to all matching environments.

Additionally, because LLM inference takes several seconds, continuous environment execution would cause \textbf{state drift} (latency-induced non-stationarity)~\parencite{Katsikopoulos_Engelbrecht_2003}, rendering the LLM's eventual decision outdated. Following established LLM-agent practices~\parencite{Zhu_etal_2023}, any environment awaiting a response is paused. This ensures the Meta-Controller's decision is applied to the exact state it observed, maintaining strict temporal alignment between strategy and execution.

% ──────────────────────────────────────────────────────────────────────────
\subsection{One-Hot Observation Augmentation}
\label{sec:onehot_aug}

Once the Commander has selected a discrete Option, that command must be physically transmitted to the worker. In the MAPPO framework, this is achieved by appending a one-hot vector directly into the agent's localized observation space, effectively acting as a ``radio receiver'' for the Commander's orders. We can see how exactly in Equation \ref{eq:hrl_obs}.

\subsubsection{Dynamic Enemy State}

Unlike the Mixing Board approach, which utilized the baseline static environment, the HRL architecture operates in an extended environment that introduces a \emph{dynamic enemy}. In this variant, the enemy is no longer a static zone interaction. Instead, the enemy tracks the agents' positions and actively chases any agent that crosses the river, creating a real-time pursuit challenge. The combat system is designed to reward coordination: if both agents attack the enemy simultaneously, they deal 50~HP of damage per step, compared to only 10~HP for a solo attack. While this 5$\times$ coordination multiplier was intended to make the \texttt{FIGHT\_ENEMY} option inherently cooperative, the overarching time penalty ($-0.01$ per step) often creates conflicting incentives. As explored in Chapter~\ref{ch:experiments}, the agents occasionally discover that having one agent rush the enemy while the other crafts is more time-efficient than waiting to launch a coordinated assault.

To give the agents spatial awareness of this dynamic threat, the three dimensions that were occupied by the goal embedding in the Mixing Board approach (indices 2--4) are repurposed to carry the enemy's real-time state: its $(x, y)$ position and its normalized health $h/100$. Providing this state information is standard practice in fully observable Markov Decision Processes (MDPs); it simply represents the agents' physical line-of-sight, which is required for the low-level PPO policy to learn the kinetic micro-actions of pursuit and combat.

\subsubsection{Observation Construction}

The vector input to the MAPPO network in this HRL mode is expanded to a 25-dimensional vector:
\begin{equation}
\label{eq:hrl_obs}
    \mathbf{v}_t^{(i)}
    =
    \left[
        \underbrace{\mathbf{p}^{(i)}}_{\mathbb{R}^2}
        \;\Big\|\;
        \underbrace{\mathbf{d}^{\text{enemy}}}_{\mathbb{R}^3}
        \;\Big\|\;
        \underbrace{\mathbf{x}}_{\mathbb{R}^{10}\text{ (inv)}}
        \;\Big\|\;
        \underbrace{\mathbf{e}_{o_t^{(i)}}}_{\mathbb{R}^{10}\text{ (one-hot)}}
    \right]
    \;\in\; \mathbb{R}^{25}
\end{equation}
where $\mathbf{p}^{(i)}$ is the agent's spatial position, $\mathbf{d}^{\text{enemy}} = (x_{\text{enemy}},\, y_{\text{enemy}},\, h_{\text{enemy}}/100)$ is the dynamic enemy state vector, $\mathbf{x}$ is the shared inventory vector, and $\mathbf{e}_{o}\in\{0,1\}^{10}$ is the one-hot encoding of the active option.

For example, if an agent is located at coordinates $(3, 4)$, the enemy is at $(5, 5)$ with 50\% health, the inventory contains exactly 1 wood, and the Commander has assigned the first option (e.g., \texttt{COLLECT\_WOOD}, index 0), the augmented vector is constructed as follows:
\begin{equation}
    \mathbf{v}_t^{(i)}
    =
    \left[
        \underbrace{3,\; 4}_{\mathbf{p}^{(i)}}
        \;\Big\|\;
        \underbrace{5,\; 5,\; 0.5}_{\mathbf{d}^{\text{enemy}}}
        \;\Big\|\;
        \underbrace{1,\; 0,\dots, 0}_{\mathbf{x} \in \mathbb{R}^{10}}
        \;\Big\|\;
        \underbrace{1,\; 0,\dots, 0}_{\mathbf{e}_o \in \mathbb{R}^{10}}
    \right]
\end{equation}

This design expands the $\mathbb{R}^{15}$ observation space of the Mixing Board approach to $\mathbb{R}^{25}$. Rather than increasing the spatial observation dimensionality, the three dimensions previously occupied by the continuous goal embedding are efficiently repurposed to carry the enemy state. The additional 10 dimensions are strictly dedicated to the one-hot option vector.

By appending this one-hot vector $\mathbf{e}_{o_t^{(i)}}$ directly to the observation space, the MAPPO network learns an \emph{option-conditional policy} $\pi_\theta(a \mid s, o)$. This one-hot encoding acts as a mathematical switch, allowing a single neural network to dynamically shift its behavior based on the Commander's active assignment, avoiding the need to train a separate policy network for each of the 10 distinct options.

% ──────────────────────────────────────────────────────────────────────────
\subsection{Bipartite Reward Structure}
\label{sec:bipartite_reward}

In standard RL, the environment reward is the ground truth. However, in our sparse environment, relying solely on this is analogous to receiving feedback only at the very end of a long episode. To ensure the agents learn the day-to-day behaviors required to achieve that final milestone, the HRL approach uses a bipartite composite reward: a sparse environment reward and a dense intrinsic reward.

\subsubsection{Extrinsic and Intrinsic Components}
\label{sec:intrinsic_reward}

As established in Section~\ref{sec:env_design}, the extrinsic environment reward $r_t^{\text{env}}$ is the true, unmodified sparse signal from the grid-world (refer to Table~\ref{tab:env_rewards}). It consists of a constant time-step penalty of $-0.01$ to encourage speed, plus massive milestone bonuses upon successfully completing a physical subtask.

Conversely, the intrinsic (shaping) reward is a dense synthetic shaping signal generated internally by the architecture to keep the agent motivated toward the Commander's assigned target. 
It consists of two sub-components: a continuous distance-based shaping mechanism, and a discrete option success bonus.

First, the continuous distance-based shaping reward toward the active option's target zone is computed identically to the PBRS framework~\parencite{Ng_Harada_Russell_1999} established in Section~\ref{subsec:pbrs}, but applied to the option's target rather than the planner's target:
\begin{equation}
\label{eq:hrl_shaping}
    r_t^{\text{shape},(i)}
    \;=\;
    \mathds{1}_{\{o^{(i)} \neq \texttt{IDLE}\}}
    \;\cdot\;
    \lambda_{\text{intr}}
    \bigl(\gamma\,\Phi^{(i)}(s_{t+1}, z_o) - \Phi^{(i)}(s_t, z_o)\bigr)
\end{equation}
where $z_o$ is the target zone of option~$o$ (\cref{tab:option_props}), $\Phi$ is the potential function, and $\lambda_{\text{intr}} = 0.05$ is the shaping scaling factor (equivalent to $\lambda_{\text{scale}}$ from the Mixing Board approach).

As established in previous chapters, this canonical difference form $\gamma\,\Phi(s_{t+1}) - \Phi(s_t)$ guarantees policy invariance and naturally prevents the agents from exploiting the reward signal by stepping back and forth~\parencite{Ng_Harada_Russell_1999}.

The indicator function $\mathds{1}_{\{o^{(i)} \neq \texttt{IDLE}\}}$ acts as a coordination mask. An agent assigned the \texttt{IDLE} option is actively waiting for its partner to satisfy a dependency; applying distance-based shaping here would incorrectly penalize it for remaining stationary. The mask zeroes out this shaping signal, allowing the agent to safely wait without destabilizing the value network.

Second, when an agent completes its active option (the relevant inventory slot increases), the distance shaping is overridden, and it receives a flat intrinsic success bonus of $+1.0$:
\begin{equation}
\label{eq:success_bonus}
    r_t^{\text{intr},(i)}
    =
    \begin{cases}
        1.0 & \text{if option } o^{(i)} \text{ terminated successfully} \\[2pt]
        r_t^{\text{shape},(i)} \text{ (\cref{eq:hrl_shaping})} & \text{otherwise}
    \end{cases}
\end{equation}

Success is detected by comparing the inventory before and after the environment step.

\subsubsection{Composite Reward}

The total reward signal fed to the PPO algorithm is a weighted sum of the environment diploma and the intrinsic breadcrumbs:
\begin{equation}
\label{eq:hrl_total_reward}
\boxed{
    r_t^{\text{total}}
    = r_t^{\text{env}}
    + \alpha_{\text{intr}} \cdot r_t^{\text{intr}}
}
\end{equation}
with $\alpha_{\text{intr}} = 0.5$. This coefficient was chosen to balance the dense intrinsic signal against the sparse environmental milestones; too large a value causes the policy to over-optimise for intrinsic shaping at the expense of actual task completion.

Note the key difference from the Mixing Board approach: here, no $\lambda$ decay schedule is applied. The intrinsic coefficient remains fixed because the LLM continuously updates the \emph{direction} (which option to pursue) rather than the magnitude of shaping.

% ──────────────────────────────────────────────────────────────────────────
\subsection{Mathematical Impact on the MAPPO Objective}
\label{sec:hrl_math_impact}

Unlike the Dynamic approach, which acted purely as an environment architect through the reward signal, the Meta-Controller strikes the MAPPO equations in \emph{two places simultaneously}: the observation space (Actor) and the reward space (Critic).

On the Actor side, the one-hot option conditioning transforms the decentralized policy of agent $i$ from $\pi_\theta(a_t^{(i)} \vert{} s_t^{(i)})$ into an option-conditional distribution $\pi_\theta(a_t^{(i)} \vert{} s_t^{(i)}, o_t^{(i)})$. 
The Actor's goal in PPO is to maximize the probability of good actions, but it uses a ``clipped objective'' to ensure the network does not update its weights too aggressively in a single step. 
By injecting the discrete option $o_t^{(i)}$, the probability ratio $r_t^{(i)}(\theta)$ between the new and old policy becomes option-aware:
\begin{equation}
    r_t^{(i)}(\theta) = \frac{\pi_\theta(a_t^{(i)} \vert{} s_t^{(i)}, o_t^{(i)})}{\pi_{\theta_{\text{old}}}(a_t^{(i)} \vert{} s_t^{(i)}, o_t^{(i)})}
\end{equation}
and the clipped surrogate objective updates accordingly:
\begin{equation}
    L^{\text{CLIP}}(\theta) = \mathbb{E} \left[ \min\left( r_t^{(i)}(\theta) A_t^{(i)},\; \text{clip}\bigl(r_t^{(i)}(\theta), 1 - \epsilon, 1 + \epsilon\bigr) A_t^{(i)} \right) \right]
\end{equation}

On the Critic side, learning is driven by the Advantage function ($A_t^{(i)}$). Remember that the Advantage measures how much better a chosen action was compared to the baseline expectation. This baseline is provided by the Centralized Value function $V_\phi$. 
To maintain the CTDE paradigm established in Section~\ref{subsec:ctde}, the Critic evaluates the \emph{global} environment state $\mathbf{s}_t$. In this HRL architecture, it also receives the joint options of all agents $\mathbf{o}_t = (o_t^{(0)}, \dots, o_t^{(N)})$. The bipartite reward feeds into this option-conditional Advantage:
\begin{equation}
    A_t^{(i)} = r_t^{\text{total}, (i)} + \gamma V_\phi(\mathbf{s}_{t+1}, \mathbf{o}_{t+1}) - V_\phi(\mathbf{s}_t, \mathbf{o}_t)
\end{equation}
Because the Value function $V_\phi$ now receives the globally augmented state $\mathbf{s}_t' = [\mathbf{s}_t \parallel \mathbf{e}_{\mathbf{o}_t}]$, the Critic learns to predict entirely different baseline expected returns depending on the combination of options actively guiding the agents.

The consequence is a tightly coupled learning loop: the Actor gradient explicitly links the local physical state $s_t^{(i)}$, the Commander's specific order $o_t^{(i)}$, and the intrinsic success of that command through $A_t^{(i)}$. The network learns distinct sub-policies for each option, not through separate networks, but through a single policy that conditions its behavior on which option the LLM has assigned.

% ──────────────────────────────────────────────────────────────────────────
\subsection{Execution Lifecycle and Stability Guardrails}
\label{sec:lifecycle_guardrails}

To ensure the HRL architecture remains robust over millions of training steps, several lifecycle management and stability mechanisms are integrated directly into the environment loop.

\subsubsection{Episode Initialization and Reset}
When an environment reaches a terminal state (gold mined, game over, or truncation at 800 steps), the environment is hard-reset. The agents are assigned initial options corresponding to the first layer of the DAG (\texttt{COLLECT\_WOOD} for the lumberjack, \texttt{COLLECT\_STONE} for the miner), ensuring they begin collecting resources immediately upon episode start without waiting for an initial LLM query.

\subsubsection{Staleness Detection and Counterfactual Reflection}
During execution, an internal watchdog tracks the age (number of steps) of each active option. If an option exceeds $T_{\text{stale}}=150$ steps without termination, the environment is marked as stale and the LLM is re-queried. The \texttt{IDLE} options are excluded from this trigger, as intentional waiting is a valid cooperative strategy.

When a staleness trigger fires, the LLM receives a \emph{counterfactual reflection prompt} instead of the standard assignment prompt. Rather than simply asking for a new assignment, this prompt forces the LLM to diagnose \textbf{why} its previous assignment failed (e.g., a missing DAG precondition like \texttt{MINE\_IRON} without a pickaxe) before issuing new options:

\begin{hegquote}
\begin{verbatim}
### REFLECTION REQUIRED:
Agent 0 was assigned "{a0_option}" and Agent 1 was assigned "{a1_option}".
After {stale_steps} steps, NEITHER has succeeded.

### YOUR TASK:
1. First, diagnose what prerequisite is MISSING.
2. Then assign corrected options that respect the DAG.

Respond ONLY with JSON:
{
  "reflection": "<diagnosis>",
  "dag_check": "<verification>",
  "agent_0_option": "<Option>",
  "agent_1_option": "<Option>"
}
\end{verbatim}
\end{hegquote}

This mechanism leverages the LLM's zero-shot reasoning to debug its own prior hallucinations, synergizing reasoning with counterfactual critique~\parencite{Yao_Zhao_Yu_Du_Shafran_Narasimhan_Cao_2023, Shinn_Cassano_Berman_Gopinath_Narasimhan_Yao_2023}.

\subsubsection{Target KL Divergence Guard}
To protect the policy network from destructive updates when intrinsic reward signals shift abruptly after an LLM reassignment, the training loop employs a target KL divergence early-stopping mechanism~\parencite{Schulman_Wolski_Dhariwal_Radford_Klimov_2017}. During each PPO epoch, the approximate KL divergence between the old and new policies is monitored:
\begin{equation}
\label{eq:approx_kl}
    D_{\text{KL}}^{\text{approx}}
    = \mathbb{E}\bigl[-\log \pi_{\theta_{\text{new}}}(a \mid s, o)
      + \log \pi_{\theta_{\text{old}}}(a \mid s, o)\bigr]
\end{equation}
If $D_{\text{KL}}^{\text{approx}} > 0.015$ (an empirical threshold widely adopted in standard PPO implementations), the remaining PPO epochs are aborted for that update, ensuring stable policy convergence.

% ──────────────────────────────────────────────────────────────────────────
\subsection{Comparison with the Mixing Board Approach}
\label{sec:hrl_vs_dyn}

Before detailing the specific fine-tuning pipeline required for the Meta-Controller, it is instructive to contrast this architecture with the Dynamic Reward Shaping (Mixing Board) approach developed in Chapter 5. While both paradigms seek to integrate LLM reasoning into MARL, they operate at fundamentally different levels of abstraction. 

The Mixing Board is a \emph{soft integration} strategy: the LLM gently nudges the priority of intrinsic rewards, but the MARL agents remain entirely responsible for discovering the physical behaviors. Conversely, the Meta-Controller represents a \emph{hard integration}: the LLM acts as a strict supervisor, dictating explicit macro-actions (options) that the agents must execute. This shift from continuous priority logits to discrete option assignments necessitates significant structural changes, particularly regarding LLM query frequency, observation augmentation, and safety mechanisms. These core architectural differences are summarized in Table~\ref{tab:approach_comparison}.

\begin{table}[H]
\centering
\caption{Structural comparison of the two LLM integration strategies.}
\label{tab:approach_comparison}
\small
\begin{tabular}{@{}p{3.5cm}p{5cm}p{5cm}@{}}
\toprule
\textbf{Aspect} & \textbf{Mixing Board (Soft Integration)} & \textbf{Meta-Controller (Hard Integration)} \\
\midrule
LLM output space &
    Continuous logits $\in \mathbb{R}^7$ &
    Discrete options $\in \mathcal{O}^2$ \\
LLM controls &
    Subtask priority weights &
    Explicit per-agent macro-actions \\
Trigger mechanism &
    Plateau + Z-score on TD error &
    Option success or staleness \\
Query granularity &
    Global (once per trigger) &
    Per-environment (batched) \\
Observation augmentation &
    None (blind to LLM) &
    One-hot option $\in \{0,1\}^{10}$ \\
Reward structure &
    PBRS with $\lambda$ decay &
    Bipartite: sparse env + dense intrinsic (fixed $\alpha$) \\
LLM query frequency &
    Low (plateau-triggered) &
    High (per-option completion) \\
Safety mechanism &
    Softmax budget, EMA, rollback &
    DAG in prompt, reflection, env freezing \\
\bottomrule
\end{tabular}
\end{table}

% ──────────────────────────────────────────────────────────────────────────
\subsection{QLoRA Fine-Tuning Pipeline}
\label{sec:qlora_pipeline}

While Section~\ref{sec:qlora} established the theoretical foundations of Quantized Low-Rank Adaptation, this section details the complete practical pipeline used to produce the specialized Commander adapter deployed throughout the HRL experiments evaluated in Chapter~\ref{ch:experiments}. As will be demonstrated, a surprisingly brief fine-tuning process on a modest synthetic dataset is sufficient to dramatically improve the Meta-Controller's decision quality.

\subsubsection{Oracle Dataset Generation}

The training dataset is generated synthetically to ensure exhaustive coverage of the crafting DAG. Exploratory successful runs from the agents helped build the dataset. A deterministic oracle systematically enumerates all meaningful inventory combinations and applies hardcoded logic to determine the optimal \texttt{(agent\_0\_option, agent\_1\_option)} assignment for each state. For example: ``if no pickaxe and wood $\ge 1$ and stone $\ge 1$, then both agents should \texttt{CRAFT\_PICKAXE}.''

To prevent the model from over-indexing on trivially simple early-game states, the generated states are grouped into six progression phases (e.g., early collection, iron mining, endgame combat). Each phase is upsampled equally to ensure a balanced training signal.

Two augmentation techniques are applied to improve generalization:

\begin{enumerate}
    \item \textbf{Stochastic Oracle Augmentation (SOA):} To prevent the model from overfitting to the perfectly logical inventory progressions generated by the script, a randomly selected inventory field is zeroed out with a 10\% probability before querying the oracle. For example, a valid mid-game state like \texttt{[iron: 1, pickaxe: 1]} might be corrupted to the physically impossible state \texttt{[iron: 1, pickaxe: 0]}. By passing this corrupted state to the oracle to generate the training label, the LLM is explicitly trained on how to gracefully recover from weird, out-of-distribution inventory combinations. This robustness is critical because observation delays and asynchronous race conditions during live RL rollouts frequently produce these mathematically inconsistent states.
    \item \textbf{Reflection Augmentation (40\% injection rate):} A synthetic hallucinated failure is prepended to the prompt: ``\emph{You previously chose [Incorrect Option]. It failed after [N] steps.}'' The model must then output the correct option despite this misleading prior. While intended to teach the Commander to recover from logical errors, Section~\ref{sec:kl_guard} demonstrates that this training signal inadvertently caused the model to assume \emph{any} reported failure implies a logical mistake (the ``Reflection Contradiction'').
\end{enumerate}

The final raw dataset contains \textbf{1,200 balanced examples} spanning all DAG phases and agent status combinations.

\subsubsection{ChatML Formatting and Dataset Splits}

The raw prompt and completion pairs are converted into Qwen's native ChatML format~\parencite{Yang_Li_Yang_Zhang_Hui_Zheng_Yu_Gao_Huang_Lv_et_al._2025}, structuring each example as a three-turn conversation (\texttt{system}, \texttt{user}, \texttt{assistant}). The system message contains the Commander's role description and DAG rules, the user message provides the current inventory state (and, for reflection examples, the failure context), and the assistant message contains the expected JSON output.

After a 90/10 train--validation split and a $2\times$ training set duplication (to ensure sufficient gradient updates given the small dataset), the final corpus comprises \textbf{2,160 training examples} and \textbf{120 validation examples}.

\subsubsection{Training Configuration}

Fine-tuning was performed on the local workstation's NVIDIA RTX 5070 GPU (12GB VRAM), as described in Section~\ref{sec:env_hardware}. The strict memory envelope forced several architectural compromises documented in Table~\ref{tab:qlora_training_config}.

\begin{table}[H]
    \centering
    \caption[QLoRA Fine-Tuning Hyperparameters]{QLoRA fine-tuning hyperparameters. Values marked with $\dagger$ were constrained by the 12GB VRAM budget.}
    \label{tab:qlora_training_config}
    \begin{tabular}{@{}llr@{}}
        \toprule
        \textbf{Category} & \textbf{Parameter} & \textbf{Value} \\
        \midrule
        Base Model & Architecture & Qwen 2.5 7B-Instruct \\
        & Source & HuggingFace Hub~\parencite{Wolf_Debut_Sanh_Chaumond_Delangue_Moi_Cistac_Rault_Louf_Funtowicz_et_al._2020} \\
        \midrule
        Quantization & Precision & 4-bit NormalFloat (NF4) \\
        & Compute dtype & bfloat16 \\
        & Double quantization & Enabled \\
        \midrule
        LoRA Adapter & Rank $r$ & 16 \\
        & Scaling $\alpha$ & 32 \\
        & Dropout & 0.05 \\
        & Target modules & All 7 (q/k/v/o/gate/up/down) \\
        \midrule
        Training & Epochs & 1 \\
        & Batch size$^\dagger$ & 1 \\
        & Gradient accumulation & 8 (eff.\ batch = 8) \\
        & Learning rate & $2 \times 10^{-4}$ \\
        & Scheduler & Cosine with warmup (10\%) \\
        & Optimizer & Paged AdamW 8-bit~\parencite{bitsandbytes} \\
        & Max sequence length$^\dagger$ & 1,024 tokens \\
        & Gradient checkpointing$^\dagger$ & Enabled \\
        \midrule
        Infrastructure & Library & HuggingFace TRL~\parencite{trl} \\
        & Adapter library & PEFT v0.19.1~\parencite{peft} \\
        & Logging & TensorBoard \\
        \bottomrule
    \end{tabular}
\end{table}

A critical training detail is the application of \emph{response-only loss masking}, implemented via the TRL library~\parencite{trl}. By default, causal language models learn by predicting every subsequent word in a sequence, meaning they would normally compute an error gradient for the prompt itself. This wastes training capacity on memorizing instructions. To solve this, a custom data collator leverages PyTorch's native \texttt{ignore\_index} mechanic: it sets the target labels of all system and user prompt tokens to $-100$. Mathematically, the Cross-Entropy loss function ignores any token with a $-100$ label, computing zero gradient for it. For example, in the sequence \texttt{[User] State is X [Assistant] \{"opt": 1\}}, every token preceding the JSON response is assigned a $-100$ label, while only the JSON characters receive valid target IDs. As a result, the model is penalized exclusively on its final JSON output, ensuring that every weight update directly improves its strategic decision-making rather than memorizing boilerplate text.

\subsubsection{Training Results and Adapter Deployment}

The complete fine-tuning process converged in approximately \textbf{3.1 hours} (1,910 steps) on the RTX 5070. As shown in Figure~\ref{fig:qlora_training}, the training loss dropped rapidly from 1.605 to 0.165, with the final evaluation metrics confirming near-perfect task acquisition:

\begin{itemize}
    \item Evaluation loss: \textbf{0.0007}
    \item Mean token accuracy: \textbf{99.98\%}
\end{itemize}

\begin{figure}[H]
    \centering
    \includegraphics[width=0.85\textwidth]{../../plots/qlora/QLoRA_Training_Performance.png}
    \caption[QLoRA Adapter Fine-Tuning Performance]{QLoRA adapter fine-tuning performance over 1,910 training steps. The training loss (left axis) converges rapidly within the first 500 steps, while the mean token accuracy (right axis) saturates near 100\%.}
    \label{fig:qlora_training}
\end{figure}

The resulting LoRA adapter weighs only \textbf{154~MB}, approximately 1\% of the full 15~GB base model, yet encodes the complete DAG reasoning and JSON formatting capabilities required by the Meta-Controller. For deployment during RL training, the adapter is loaded via the PEFT library and merged with the base model in memory, or alternatively converted to GGUF format and served through the Ollama inference framework (Section~\ref{sec:env_hardware}). Both pathways produce inference latencies compatible with the synchronous RL training loop, as quantified in Section~\ref{sec:computational_overhead}. This fine-tuned adapter serves as the Meta-Controller for all ``QLoRA'' HRL configurations evaluated in Chapter~\ref{ch:experiments}.

% ──────────────────────────────────────────────────────────────────────────
\subsection{Summary}
\label{sec:hrl_summary}

The Meta-Controller approach establishes a strict separation of concerns: the LLM Commander dictates \emph{what to do} (strategic option selection), while the MAPPO policy learns \emph{how to do it} (kinetic execution). By appending a one-hot option encoding to the state vector, a single policy network learns option-specific behaviors. Simultaneously, the bipartite reward structure ensures the policy receives both dense intrinsic guidance and sparse environmental grounding.

Crucially, this architecture functions as an automated curriculum generator. Unlike traditional Curriculum Learning (Section~\ref{sec:curriculum_learning}), which requires manually designing environments of escalating difficulty, the Meta-Controller orchestrates a skill curriculum dynamically. It guides agents from basic resource gathering to advanced combat coordination entirely within a single, static environment.

\begin{table}[H]
\centering
\caption{Key hyperparameters of the Meta-Controller approach.}
\label{tab:hrl_hyperparams}
\begin{tabular}{@{}lll@{}}
\toprule
\textbf{Component} & \textbf{Parameter} & \textbf{Value} \\
\midrule
Options          & $|\mathcal{O}|$                & 10 \\
Staleness        & $T_{\text{stale}}$             & 150 steps \\
Cooldown         & Post-query cooldown            & 50 steps \\
Intrinsic reward & Distance scale $\lambda_{\text{intr}}$ & 0.05 \\
                 & Success bonus                  & $+1.0$ \\
Composite reward & $\alpha_{\text{intr}}$         & 0.5 \\
                 & $D_{\max}$                     & 85.0 \\
Target KL        & KL threshold                   & 0.015 \\
Network          & One-hot dimension              & 10 \\
                 & Total \texttt{vec\_dim}        & 25 \\
\bottomrule
\end{tabular}
\end{table}

\section{Empirical Results: Breaking the Exploration Bottleneck}
\label{ch:experiments}

This chapter presents the experimental setup, the environments used for evaluation, and the results obtained. We specifically focus on validating the hypothesis that a Meta-Controller LLM can effectively resolve exploration bottlenecks in a sparse Multi-Agent Reinforcement Learning (MARL) environment through dynamic reward shaping, and subsequently, how Hierarchical Reinforcement Learning (HRL) further improves these capabilities.

\subsection{Experimental Setup}
\label{sec:exp_setup}

To rigorously evaluate the proposed architectures, we compare distinct configurations within our multi-agent resource-gathering environment, structured into two parts. 

For the \textbf{Dynamic Reward Shaping} experiments, we evaluate:
\begin{enumerate}
    \item \textbf{Sparse Baseline:} A standard MAPPO implementation~\parencite{Yu_Velu_Vinitsky_Gao_Wang_Bayen_Wu_2022} relying purely on sparse environmental rewards (+1 for completing milestones).
    \item \textbf{LLM Dynamic (Gemini 2.5 Flash Lite):} The proposed ``Mixing Board'' architecture, utilizing Google's Gemini 2.5 Flash Lite model via API as the Commander to dynamically shift potential-based reward weights. 
    \item \textbf{LLM Dynamic (Local Qwen 2.5 7B):} An equivalent dynamic architecture utilizing a localized, smaller parameter model (Qwen 2.5 7B) to observe the efficacy of smaller, self-hosted LLMs. 
\end{enumerate}

For the \textbf{Hierarchical Reinforcement Learning (HRL)} experiments, we evaluate three configurations:
\begin{enumerate}
    \item \textbf{No-LoRA Base (Qwen 2.5 7B):} The un-finetuned base model serving as the HRL Meta-Controller (the Commander).
    \item \textbf{QLoRA with Reflection (Qwen 2.5 7B):} A QLoRA fine-tuned model (as detailed in Section~\ref{sec:qlora_pipeline}, the Qwen 2.5 7B model was fine-tuned using a 4-bit NF4 quantization pipeline on 1,200 synthetic oracle examples) with counterfactual reflection triggered on option staleness.
    \item \textbf{QLoRA without Reflection (Qwen 2.5 7B):} The same fine-tuned model, with the reflection mechanism disabled.
\end{enumerate}

All runs were executed with identical MAPPO hyperparameters (clipping $\epsilon=0.2$, GAE\footnote{Generalized Advantage Estimation} $\lambda=0.95$, $\gamma=0.99$, across 128 parallel environments).

% ==========================================
% PART 1: DYNAMIC REWARD SHAPING
% ==========================================
\subsection{Part 1: Dynamic Reward Shaping}
\label{sec:exp_dynamic}

The most crucial metric validating the LLM Dynamic architecture is the progression through the environment's strict dependency Directed Acyclic Graph (DAG). 

\begin{figure}[H]
    \centering
    \includegraphics[width=\textwidth]{../../plots/baseline_vs_llmdynamic_results/subtasks_comparison.png}
    \caption{Subtask Completion Rates: Comparison between the Sparse Baseline and the LLM Dynamic models over 2000 training updates.}
    \label{fig:subtasks_comparison_dynamic}
\end{figure}

As shown in \cref{fig:subtasks_comparison_dynamic}, all three configurations conquer the early-game tasks (Wood, Stone) almost immediately. The Pickaxe subtask follows shortly after, with the Baseline reaching it at step~1.1M and Gemini at step~720K.

However, the Baseline completely collapses at the \emph{Iron} bottleneck. Without any dynamic reward shaping to guide exploration, the Baseline workers never discover how to consistently mine Iron (0\% final rate), resulting in a total failure to learn all downstream tasks (Sword, Armor, Bridge). 

In contrast, both LLM-guided runs break through the Iron bottleneck, but at remarkably different speeds. \textbf{Qwen} unlocks Iron at step~10M, subsequently discovering Sword and Armor at essentially the same step (10.1M). This rapid cascade demonstrates the snowball effect: once a critical bottleneck is broken, downstream tasks are unlocked quickly. \textbf{Gemini 2.5 Flash Lite} unlocks Iron much later, at step~40.7M (approximately 4$\times$ slower than Qwen). However, once unlocked, the same cascade occurs: Sword and Armor appear at step~41.3M.

Neither run achieved Enemy defeats (0\% for all three configurations), indicating that the combat subtask represents a qualitatively different challenge beyond reward gravity, likely requiring specialised action coordination that the current shaping signal alone cannot provide.

\subsubsection{LLM Interventions in Action}
\label{sec:exp_interventions}

To understand exactly \emph{how} the LLMs broke the bottleneck, we analyse the evolution of the reward weights. The weights shown in \cref{fig:llm_weights_evolution} are the \emph{post-softmax} weights actually applied to the potential-based shaping.

\begin{figure}[H]
    \centering
    \begin{subfigure}{0.8\textwidth}
        \includegraphics[width=\textwidth]{../../plots/baseline_vs_llmdynamic_results/llm_weights_evolution_gemini.png}
        \caption{Gemini 2.5 Flash Lite Meta-Controller (198 interventions)}
    \end{subfigure}
    \vskip\baselineskip
    \begin{subfigure}{0.8\textwidth}
        \includegraphics[width=\textwidth]{../../plots/baseline_vs_llmdynamic_results/llm_weights_evolution_qwen.png}
        \caption{Local Qwen 2.5 7B Meta-Controller (20 interventions)}
    \end{subfigure}
    \caption{Evolution of the Dynamic Reward Weights (post-softmax, as applied to the PBRS signal).}
    \label{fig:llm_weights_evolution}
\end{figure}

\Cref{fig:llm_weights_evolution} reveals strikingly different command strategies. \textbf{Gemini's strategy} as a Commander is characterised by \emph{broad amplification} of the early-chain tasks. $w_{\text{wood}}$ and $w_{\text{stone}}$ rapidly climb and stay high, while $w_{\text{iron}}$ remains near zero after the softmax (mean~0.01). This reveals a counter-intuitive insight: despite the raw LLM outputs assigning $w_{\text{iron}} = 1.0$, the \emph{softmax normalisation} compresses this into a negligible effective weight because the Commander simultaneously assigns high values to wood, stone, and bridge.

\textbf{Qwen's strategy} exhibits a \emph{focused suppression-and-amplification} profile. After an initial exploration phase, $w_{\text{wood}}$ and $w_{\text{stone}}$ are aggressively dropped, while $w_{\text{iron}}$ rises sharply to a peak of 3.87 during the mid-training bottleneck. By zeroing out already-mastered tasks and concentrating the magnetic pull entirely on the bottleneck, Qwen breaks through Iron 30M steps earlier than Gemini.

\subsubsection{Reward Analysis}
\label{sec:exp_rewards}

\begin{figure}[H]
    \centering
    \includegraphics[width=0.7\textwidth]{../../plots/baseline_vs_llmdynamic_results/reward_comparison.png}
    \caption{Average Environment Reward (True Milestones Reached).}
    \label{fig:env_rewards_dynamic}
\end{figure}

\Cref{fig:env_rewards_dynamic} presents the true average environment reward. Both LLM Dynamic runs achieve significantly higher returns than the Baseline, reaching an identical peak reward ($\approx$0.076). This suggests that once the bottleneck is broken, the final quality of the learned policy is comparable; the difference lies purely in the \emph{sample efficiency} driven by Qwen's highly focused interventions.

However, a closer inspection of the training trajectories reveals a critical difference in policy stability and retention. Around step $2.97 \times 10^7$, the Qwen-guided workers exhibit a sharp drop in performance, seemingly forgetting how to complete the Bridge subtask. This catastrophic forgetting is reflected in a corresponding dip in Qwen's average environment reward. Conversely, the Gemini Commander demonstrates superior stability: once it teaches the workers to build a Bridge, they retain this behavior. This retention culminates in a higher average reward spike for Gemini around step $4.65 \times 10^7$. These observations suggest that while Qwen's aggressive, narrowly focused shaping breaks bottlenecks much faster, Gemini's broader amplification strategy (\Cref{fig:llm_weights_evolution}) may foster a more robust, generalizable policy that is less susceptible to catastrophic forgetting.

% ==========================================
% PART 2: HIERARCHICAL REINFORCEMENT LEARNING
% ==========================================
\subsection{Part 2: Hierarchical Reinforcement Learning}
\label{sec:exp_hrl}

While the dynamic shaping architecture successfully guided the lumberjack and miner past initial bottlenecks (like Iron), it ultimately stalled at the combat phase (Enemy, Gold). Modifying the continuous reward landscape was insufficient to teach the workers the discrete, long-horizon macro-actions required for coordinated combat.

This limitation motivated the shift to a full \textbf{Hierarchical Reinforcement Learning (HRL)} architecture, where the Commander does not merely shape the rewards, but explicitly radios discrete \emph{Options} (e.g., ``Craft Sword'', ``Defeat Enemy'') directly to the lower-level MAPPO policy. Three configurations were evaluated over 2000 optimization updates ($\approx$65.5M environment steps each):

\begin{enumerate}
    \item \textbf{No-LoRA Base:} The un-finetuned Qwen 2.5 7B model serving as Meta-Controller.
    \item \textbf{QLoRA (With Reflection):} The QLoRA fine-tuned model with counterfactual reflection logic, triggered when an assigned Option becomes stale (workers fail to complete it within 150 steps).
    \item \textbf{QLoRA (No Reflection):} The identical fine-tuned model, but with the reflection trigger suppressed. On staleness, the Commander re-evaluates the raw environment state without injecting any prior failure context.
\end{enumerate}

\subsubsection{Subtask Progression: Solving the DAG}

\Cref{tab:hrl_breakthrough} presents the exact environment step at which each subtask was first completed across the three configurations. Each one managed to complete the DAG at some point. 

\begin{table}[H]
\centering
\caption{First breakthrough step (millions of environment steps) for each subtask. \textbf{Bold} indicates the fastest configuration per subtask.}
\label{tab:hrl_breakthrough}
\begin{tabular}{lccc}
\hline
\textbf{Subtask} & \textbf{No-LoRA} & \textbf{QLoRA + Refl.} & \textbf{QLoRA -- No Refl.} \\
\hline
Wood       & \textbf{0.16M} & 0.29M & 0.26M \\
Stone      & \textbf{0.16M} & 0.29M & 0.26M \\
Pickaxe    & 13.50M & \textbf{1.25M} & 0.72M \\
Bridge     & 30.61M & \textbf{2.03M} & 1.18M \\
Iron       & 16.61M & \textbf{2.13M} & 3.01M \\
Sword      & 16.71M & \textbf{2.29M} & 3.15M \\
Armor      & 16.78M & \textbf{2.29M} & 3.15M \\
Enemy      & 30.74M & \textbf{2.49M} & 3.24M \\
Gold       & 32.31M & 7.08M & \textbf{6.29M} \\
\hline
\end{tabular}
\end{table}

Several patterns emerge from \Cref{tab:hrl_breakthrough}. First, all three configurations solve the early-game tasks (Wood, Stone) within the first 0.3M steps, with the No-LoRA base marginally faster due to lower inference overhead. Second, the QLoRA models exhibit a dramatic advantage starting at Pickaxe: both fine-tuned runs complete the mid-game milestones (Pickaxe through Armor) approximately \textbf{8--14$\times$ faster} than the base model. The No-LoRA run does not reach Pickaxe until 13.5M steps, nearly $11\times$ slower than QLoRA with Reflection (1.25M), demonstrating that fine-tuning is essential for coherent option sequencing.

Third, the No-LoRA run exhibits a pronounced \emph{ordering anomaly}: it completes Iron, Sword, and Armor ($\sim$16.7M) well before Bridge (30.6M). Since the bridge is a prerequisite for reaching the enemy side of the river, these combat-preparatory milestones accumulate without immediate use. The workers spend $\sim$14M steps fully equipped but unable to cross the river. Both QLoRA runs, by contrast, follow the optimal DAG ordering (building the bridge \emph{before} crafting combat equipment), resulting in a tight cascade from Bridge to Enemy ($<$0.5M steps gap).

Finally, at the Gold milestone, QLoRA without Reflection is the fastest (6.29M), followed closely by QLoRA with Reflection (7.08M). The No-LoRA base reaches Gold at 32.31M, \textbf{5.1$\times$ slower}, confirming that without fine-tuning, the Commander struggles to plan the late-game transition from combat to resource extraction.

\begin{figure}[H]
    \centering
    \begin{subfigure}{0.32\textwidth}
        \includegraphics[width=\textwidth]{../../plots/hrl_reflection_ablation/hrl_ablation_subtask_Wood.png}
        \caption{Wood}
    \end{subfigure}
    \hfill
    \begin{subfigure}{0.32\textwidth}
        \includegraphics[width=\textwidth]{../../plots/hrl_reflection_ablation/hrl_ablation_subtask_Stone.png}
        \caption{Stone}
    \end{subfigure}
    \hfill
    \begin{subfigure}{0.32\textwidth}
        \includegraphics[width=\textwidth]{../../plots/hrl_reflection_ablation/hrl_ablation_subtask_Pickaxe.png}
        \caption{Pickaxe}
    \end{subfigure}
    \vskip 0.5\baselineskip
    \begin{subfigure}{0.32\textwidth}
        \includegraphics[width=\textwidth]{../../plots/hrl_reflection_ablation/hrl_ablation_subtask_Iron.png}
        \caption{Iron}
    \end{subfigure}
    \hfill
    \begin{subfigure}{0.32\textwidth}
        \includegraphics[width=\textwidth]{../../plots/hrl_reflection_ablation/hrl_ablation_subtask_Sword.png}
        \caption{Sword}
    \end{subfigure}
    \hfill
    \begin{subfigure}{0.32\textwidth}
        \includegraphics[width=\textwidth]{../../plots/hrl_reflection_ablation/hrl_ablation_subtask_Armor.png}
        \caption{Armor}
    \end{subfigure}
    \vskip 0.5\baselineskip
    \begin{subfigure}{0.32\textwidth}
        \includegraphics[width=\textwidth]{../../plots/hrl_reflection_ablation/hrl_ablation_subtask_Bridge.png}
        \caption{Bridge}
    \end{subfigure}
    \hfill
    \begin{subfigure}{0.32\textwidth}
        \includegraphics[width=\textwidth]{../../plots/hrl_reflection_ablation/hrl_ablation_subtask_Enemy.png}
        \caption{Enemy}
    \end{subfigure}
    \hfill
    \begin{subfigure}{0.32\textwidth}
        \includegraphics[width=\textwidth]{../../plots/hrl_reflection_ablation/hrl_ablation_subtask_Gold.png}
        \caption{Gold}
    \end{subfigure}
    \caption{Subtask completion rates across all nine environment milestones for the three HRL configurations. The dependency chain flows left-to-right, top-to-bottom.}
    \label{fig:hrl_all_subtasks}
\end{figure}

Crucially, \Cref{fig:hrl_all_subtasks} reveals that \textbf{all three configurations ultimately solve the complete DAG}, reaching 100\% completion on every subtask including Gold. This is a qualitative leap over the Dynamic Shaping architecture (Part~1), which never reached Enemy or Gold. The key differentiator between configurations is not whether they succeed, but \emph{how quickly}: QLoRA with Reflection reaches 100\% on Gold by $\sim$10M steps, QLoRA without Reflection by $\sim$8M steps, while the No-LoRA base requires $\sim$35M steps, consuming more than half the training budget before the workers first collect Gold.

Interestingly, the QLoRA with Reflection run's subtask curves exhibit a distinctive ``staircase'' pattern (\Cref{fig:hrl_all_subtasks}, Bridge and Enemy panels): early plateaus near 0\% are followed by abrupt jumps to 100\%. This indicates that the rapid option cycling caused by the reflection mechanism prevents workers from making incremental progress; instead, the completion rate stays near zero until the PPO policy forces a breakthrough, after which it saturates instantly. By contrast, the QLoRA without Reflection run shows a much smoother, gradual learning ramp (taking roughly 4$\times$ as many updates to transition from 0\% to 90\% on the Enemy subtask compared to No-LoRA). This smooth monotonic ramp is evidence of deliberate, stable option sequencing by the fine-tuned Meta-Controller. The No-LoRA base falls in between, showing delayed but relatively fast transitions once it eventually stumbles upon the correct behavior.

\subsubsection{Reward Trajectories}

\begin{figure}[H]
    \centering
    \includegraphics[width=\textwidth]{../../plots/hrl_reflection_ablation/hrl_ablation_rewards.png}
    \caption{Average Extrinsic Reward across HRL configurations. All three runs converge to comparable final rewards ($\sim$0.326), but the QLoRA runs reach this level $\sim$3$\times$ faster than the No-LoRA base.}
    \label{fig:hrl_rewards}
\end{figure}

The reward trajectories (\Cref{fig:hrl_rewards}) quantify the sample efficiency gap. All three configurations converge to nearly identical final rewards (No-LoRA: 0.326, QLoRA+Refl: 0.326, QLoRA--NoRefl: 0.327), with peak rewards clustered at $\sim$0.335. This convergence confirms that the underlying PPO policy has sufficient capacity to master the environment regardless of the Meta-Controller quality, given enough time. The practical difference is entirely one of sample efficiency.

Compared to the Dynamic Shaping architectures (Part~1), which plateaued at 0.076, the HRL architecture achieves a \textbf{4.4$\times$ improvement} in peak reward. The QLoRA runs reach their peak reward plateau by $\sim$28M steps (No Reflection) and $\sim$37M steps (With Reflection), while the No-LoRA base does not stabilize until $\sim$55M steps. This $\sim$2$\times$ sample efficiency gap between the QLoRA runs is the first quantitative signal that the reflection mechanism, while not preventing convergence, actively delays it.

\subsubsection{Stabilizing Non-Stationary Objectives via KL-Divergence Guarding}
\label{sec:kl_guard}

In the HRL architecture, the transition between discrete macro-actions (Options) induces abrupt shifts in the intrinsic reward landscape. From the perspective of the low-level PPO workers, the objective function is highly non-stationary. To prevent catastrophic forgetting of previously learned kinetic policies, we bound the surrogate objective update using a strict Kullback-Leibler (KL) divergence guard.

During the optimization epochs, we continuously monitor the approximate KL divergence between the behavioral policy $\pi_{\theta_{\text{old}}}$ and the updating policy $\pi_{\theta}$~\parencite{Schulman_Wolski_Dhariwal_Radford_Klimov_2017}:
\begin{equation}
    D_{\text{KL}}^{\text{approx}} = \mathbb{E}\left[ -\log \left( \frac{\pi_{\theta}(a|s, o)}{\pi_{\theta_{\text{old}}}(a|s, o)} \right) \right]
\end{equation}

If $D_{\text{KL}}^{\text{approx}}$ exceeds the predefined trust region threshold ($\delta = 0.015$, implemented as \texttt{target\_kl}), the gradient updates for the current batch are immediately aborted.

\begin{figure}[H]
    \centering
    \includegraphics[width=\textwidth]{../../plots/hrl_reflection_ablation/hrl_ablation_kl_spikes.png}
    \caption{Maximum approximate KL divergence per update across the three HRL configurations. The dashed red line marks the early-stopping threshold $\delta = 0.015$.}
    \label{fig:kl_spikes}
\end{figure}

\begin{table}[H]
\centering
\caption{KL divergence statistics and policy stability metrics across HRL configurations.}
\label{tab:kl_stats}
\begin{tabular}{lccc}
\hline
\textbf{Metric} & \textbf{No-LoRA} & \textbf{QLoRA + Refl.} & \textbf{QLoRA -- No Refl.} \\
\hline
Mean $\max D_{\text{KL}}$       & 0.0066 & 0.0082 & 0.0073 \\
Peak $\max D_{\text{KL}}$       & 0.025  & \textbf{0.117}  & 0.107 \\
KL early-stop triggers           & 46 / 2{,}000 & \textbf{234 / 2{,}000} & 152 / 2{,}000 \\
Staleness resets (total)         & 28{,}927 & \textbf{88{,}711} & 39{,}137 \\
LLM queries dispatched           & 12{,}752 & 12{,}791 & \textbf{13{,}710} \\
Final TD error (last 100 updates)& 0.042  & 0.024  & \textbf{0.023} \\
Mean TD error (all updates)      & 0.046  & 0.060  & 0.048 \\
\hline
\end{tabular}
\end{table}

\Cref{fig:kl_spikes} and \Cref{tab:kl_stats} reveal a counter-intuitive relationship between KL instability and learning dynamics. The QLoRA with Reflection run triggers the KL early-stopping guardrail \textbf{5.1$\times$ more often} than the No-LoRA base (234 vs.\ 46 triggers) and \textbf{1.5$\times$ more often} than the No-Reflection run (234 vs.\ 152). Both QLoRA runs exhibit large peak KL spikes ($>$0.1), indicating aggressive policy shifts during option transitions. The No-LoRA base, by contrast, shows modest peak KL (0.025) and infrequent triggers. Its policy shifts are small because the un-finetuned Commander generates less decisive option switches. 

These aggressive policy shifts and high KL guard triggers in the QLoRA with Reflection run also directly manifest as increased fluctuations in its average extrinsic reward trajectory (\Cref{fig:hrl_rewards}). The constant disruption of the objective function prevents the PPO workers from smoothly optimizing their kinetic behaviors.

Critically, the QLoRA with Reflection run accumulates \textbf{2.3$\times$ more staleness resets} than the No-Reflection run (88{,}711 vs.\ 39{,}137), despite dispatching nearly the same number of LLM queries (12{,}791 vs.\ 13{,}710). This means each reflection-triggered re-query produces an Option that \emph{fails faster}---the Commander actively generates worse plans when told its previous plan failed. The rapid cycling of short-lived Options produces a high frequency of small policy shifts: each individual shift is small enough to avoid the KL guardrail, but the cumulative effect prevents the Critic from converging on a stable value estimate. We term this pattern \emph{death by a thousand cuts}: individually invisible perturbations that collectively devastate sample efficiency.

The No-Reflection run follows the opposite pattern: \emph{infrequent, large, guarded} transitions. The Commander commits to long-horizon Options and rarely switches. When it does (for example, transitioning from \texttt{CRAFT\_SWORD} to \texttt{DEFEAT\_ENEMY}), the policy shift is dramatic (peak KL = 0.107, or $7\times$ the threshold). The guardrail catches these large transitions and aborts the gradient update. Over subsequent updates, the policy absorbs the new objective gradually. This pattern of rare but controlled disruptions produces the lowest mean TD error (0.048) among the three runs.

\subsubsection{The Paradox of Reflection}
\label{sec:reflection_ablation}

The staleness data (\Cref{tab:kl_stats}) provides the mechanistic explanation for the delay caused by reflection. The With-Reflection run dispatches essentially the same number of LLM queries as the other runs ($\sim$12{,}800), yet accumulates \textbf{3.1$\times$ more staleness resets} than the No-LoRA base (88{,}711 vs.\ 28{,}927) and \textbf{2.3$\times$ more} than the No-Reflection run (88{,}711 vs.\ 39{,}137). Each reflection-triggered query thus produces a worse Option, an Option that expires faster and triggers yet another staleness event.

The root cause lies in how the instruction-tuning dataset was constructed. During QLoRA fine-tuning, the dataset synthetically injected counterfactual failures: the training prompt stated \emph{``You previously chose [Incorrect Option]. It failed.''} The model then learned to output the correct Option. Consequently, the Commander associated the reflection trigger with a strong prior: \textbf{my previous choice was logically wrong, and I must switch to a different option.}

During inference, however, option staleness is frequently triggered not by a logical error by the Commander, but by kinetic delays on the ground: the lumberjack navigating around the river terrain, the miner queuing at resource nodes, or simply needing more steps to execute a fundamentally correct plan. When the orchestrator fires the reflection prompt, the Commander, conditioned by its training data, assumes it made a reasoning error. It abandons the correct long-horizon option and switches to an alternative, derailing the workers' execution.

To illustrate this transparency, consider a snapshot directly from the evaluation logs where the workers have already defeated the Enemy and just need to collect the Gold. The Commander correctly assigns \texttt{COLLECT\_GOLD}, but when the workers take too long and trigger a staleness reset, the LLM hallucinates a logical error and regresses to an unnecessary prerequisite:

\begin{verbatim}
Inventory: {..., 'pickaxe': 1, 'sword': 1, 'bridge': 1, 'enemy': 1, 'gold': 0}
Previous Assignment: Agent 1 failed on COLLECT_GOLD

LLM Response:
{
  "reflection": "Previous assignment of IDLE and COLLECT_GOLD failed because 
                 agents lacked the necessary sub-tier items.",
  "dag_check": "A0 waits (bridge done), A1 crafts armor.",
  "agent_0_option": "IDLE",
  "agent_1_option": "CRAFT_ARMOR"
}
\end{verbatim}
Even though Armor is entirely irrelevant once the Enemy is defeated, the Commander doubts its own correct logic because the reflection prompt implied it made a mistake.

We term this failure mode the \emph{Reflection Contradiction}: the architecture inadvertently trained the LLM to doubt its own correct logic. Disabling the reflection mechanism allows the Commander to maintain stable plans. On re-query, it re-confirms its current Option rather than compulsively switching.

\begin{figure}[H]
    \centering
    \includegraphics[width=\textwidth]{../../plots/hrl_reflection_ablation/hrl_ablation_td_error.png}
    \caption{Temporal Difference Error Variance. The With-Reflection configuration exhibits persistent high variance during mid-training from contradictory option switches. Both QLoRA runs converge to similar final TD errors ($\sim$0.024), while the No-LoRA base settles higher at 0.042.}
    \label{fig:hrl_td_error}
\end{figure}

The downstream effect is quantified in the TD Error Variance (\Cref{fig:hrl_td_error}). The With-Reflection configuration maintains the highest mean TD error across the full training run (0.060), 25\% higher than the No-Reflection run (0.048) and 30\% higher than the No-LoRA base (0.046). This elevated TD error during mid-training reflects the Critic's inability to track the rapidly switching intrinsic reward landscape caused by reflection-induced option churn. By the final 100 updates, both QLoRA runs converge to comparable TD errors (With Reflection: 0.024, No Reflection: 0.023), indicating that the Critic eventually stabilizes once the policy has saturated all subtasks and option switching diminishes. The No-LoRA base, despite its lower mid-training variance, settles at a higher final TD error (0.042). Its Critic struggles with the less structured option transitions produced by the un-finetuned model.

Ultimately, all three HRL configurations solve the complete environment DAG, representing a qualitative breakthrough over the Dynamic Shaping architecture. The key finding is that \textbf{QLoRA fine-tuning is essential for sample efficiency} (5$\times$ faster to Gold than the base model), while the reflection mechanism (despite its theoretical appeal) \textbf{delays convergence} by injecting ungrounded self-doubt into the Meta-Controller's reasoning. The KL-divergence guardrail proves indispensable for absorbing the inherent non-stationarity of option transitions, but it cannot protect against the high-frequency, sub-threshold perturbations generated by the reflection loop.

% ==========================================================================
\subsection{Part 3: Robustness and Semantic Adaptation}
\label{sec:robustness_and_adaptation}

Having established that hard integration (HRL) outperforms soft integration (Dynamic Shaping) in solving the complete dependency graph, we now return briefly to the Mixing Board architecture. To fully validate the theoretical claims made in \Cref{sec:ema_smoothing} regarding stability, we isolate two critical variables of the soft integration approach: the numerical stability provided by temporal smoothing, and the zero-shot semantic adaptation capabilities of the LLM Commander.

\subsubsection{Ablation: The Necessity of EMA Smoothing}
\label{sec:ema_ablation}

In \Cref{sec:ema_smoothing}, we posited that feeding raw, discrete priority logits directly into the potential-based shaping function would induce severe non-stationarity. To empirically verify this, an ablation run was conducted wherein the Exponential Moving Average (EMA) layer was disabled.

\begin{figure}[H]
    \centering
    \includegraphics[width=0.48\textwidth]{../../plots/llm_dynamic_ablations_and_shifts/ablation_rewards.png}
    \hfill
    \includegraphics[width=0.48\textwidth]{../../plots/llm_dynamic_ablations_and_shifts/ablation_td_error.png}
    \caption{Comparison of training stability with and without EMA smoothing ($\alpha=0.2$). Disabling smoothing induces a catastrophic spike in Temporal Difference (TD) error variance (right).}
    \label{fig:ema_ablation_td}
\end{figure}

\begin{figure}[H]
    \centering
    \includegraphics[width=0.48\textwidth]{../../plots/llm_dynamic_ablations_and_shifts/ablation_subtask_Iron.png}
    \hfill
    \includegraphics[width=0.48\textwidth]{../../plots/llm_dynamic_ablations_and_shifts/ablation_subtask_Sword.png}
    \caption{Subtask completion rates for Iron (left) and Sword (right). The variance spike directly correlates with suboptimal learning and temporary regressions in advanced subtasks.}
    \label{fig:ema_ablation_subtasks}
\end{figure}

As shown in Figure~\ref{fig:ema_ablation_td}, removing the EMA filter results in a noticeable spike in TD-error variance. The ablation itself was executed by passing a \texttt{skip\_ema=True} flag to the \texttt{LLMOrchestrator}, which bypassed the exponential smoothing step entirely.

The Critic network struggles to track a reward objective that shifts instantaneously. Crucially, Figure~\ref{fig:ema_ablation_subtasks} illustrates the behavioral consequence: the non-stationarity leads to suboptimal optimization, causing the workers to exhibit slight regressions in their \texttt{Iron} and \texttt{Sword} completion rates during the mid-stages of training. However, it is worth noting that the No-EMA workers do eventually recover and climb back to a 100\% success rate on these tasks. Thus, while the workers can mathematically overcome the noise, this proves that the EMA smoothing equation is highly helpful for bridging discrete LLM outputs with continuous RL gradients, enabling faster convergence and preventing such regressions.

To quantify the optimization efficiency provided by the EMA filter, Table~\ref{tab:ema_convergence_speed} details the number of environment steps required to reliably learn (i.e., achieve a $\ge 90\%$ success rate) the advanced subtasks.

\begin{table}[H]
    \centering
    \begin{tabular}{lrr}
        \toprule
        \textbf{Subtask} & \textbf{With-EMA (Baseline)} & \textbf{No-EMA (Ablation)} \\
        \midrule
        \texttt{Wood} & 0.33M steps & 2.36M steps ($7.2\times$ slower) \\
        \texttt{Stone} & 0.82M steps & 2.46M steps ($3.0\times$ slower) \\
        \texttt{Pickaxe} & 1.31M steps & 6.88M steps ($5.3\times$ slower) \\
        \texttt{Bridge} & 2.06M steps & 9.24M steps ($4.5\times$ slower) \\
        \texttt{Iron} & 2.22M steps & 21.00M steps ($9.4\times$ slower) \\
        \texttt{Sword} & 2.52M steps & 21.63M steps ($8.6\times$ slower) \\
        \texttt{Armor} & 2.52M steps & 21.63M steps ($8.6\times$ slower) \\
        \texttt{Enemy} & 3.01M steps & Failed to Converge \\
        \texttt{Gold} & 7.07M steps & Failed to Converge \\
        \bottomrule
    \end{tabular}
    \caption{Environment steps required to reach a 90\% completion rate. Disabling the EMA smoothing filter causes a severe learning slowdown across all tasks, culminating in an inability to converge on the terminal tasks within the training envelope.}
    \label{tab:ema_convergence_speed}
\end{table}

\begin{figure}[H]
    \centering
    \includegraphics[width=0.8\textwidth]{../../plots/llm_dynamic_ablations_and_shifts/ablation_all_subtasks_bar.png}
    \caption{Maximum subtask completion rates achieved across all runs. Despite the early variance spikes, No-EMA eventually masters all tasks identical to the baseline.}
    \label{fig:ema_ablation_all_subtasks}
\end{figure}

\subsubsection{Computational, Financial, and Energy Overheads}
\label{sec:computational_overhead}

A critical practical consideration when integrating LLMs tightly into the RL training loop is the wall-clock time overhead. Because the meta-controller queries the LLM dynamically throughout the episode, relying on external or cloud-based APIs introduces severe network latency and rate-limiting bottlenecks. 

Analyzing the total execution times for a complete 128-update training run reveals this stark reality. The \textbf{Standard HRL Baseline} (which uses a traditional static meta-controller without LLM queries) completed training in \textbf{8.30 hours}. When deploying an external cloud model (Gemini 2.5 Flash-Lite) for the LLM Dynamic run, the training time ballooned to \textbf{12.27 hours}---a roughly 48\% increase purely due to API latency, despite querying a highly optimized "lite" model. Using a heavier state-of-the-art online model (such as GPT-4) was entirely infeasible for our full training comparisons due to strict API rate limits and prohibitive latency over thousands of calls.

In contrast, our optimized \textbf{LLM Dynamic (LoRA)} implementation, running a locally fine-tuned Qwen2.5-7B model, completed the identical training run in just \textbf{7.78 hours}. By keeping inference entirely local, batching requests, and avoiding network overhead, the LLM-guided approach achieved execution speeds completely on par with (and in this instance, marginally faster than) the standard HRL baseline. This proves that dynamic LLM orchestration is not just theoretically advantageous, but computationally viable when deployed efficiently at the edge.

Beyond pure execution speed, the financial and energy footprints of these architectures are critical factors for independent researchers. Relying on commercial online APIs for the meta-controller generates a staggering volume of requests; simulating 128 parallel environments over millions of steps yields millions of input and output tokens per training run. Relying on paid APIs for this volume of queries would render the approach financially prohibitive. Furthermore, reinforcement learning is highly hardware-intensive. A standard local workstation drawing approximately 750W continuously over a 12.27-hour API-bound run---or a 13.5-hour unoptimized local run---results in substantial electricity consumption, which becomes especially burdensome when multiplied by dozens of iterations, failed experiments, and hyperparameter sweeps. In regions with high energy costs such as Switzerland, the ability of the optimized local LoRA implementation to slash training time down to 7.78 hours provides a compounded benefit: it completely eliminates API token costs while simultaneously reducing the electrical footprint and associated utility costs by nearly 40\% per run.

\subsubsection{Environment Shift: The Paradox of Semantic Over-Confidence}
\label{sec:semantic_adaptation}

A fundamental critique of using LLMs in MARL is whether they rely on genuine semantic world-models or merely overfit to node names. To test this, the environment was silently modified at Update 500: the `Workbench` zone was renamed to `Furnace` in the Commander's prompt space. A secondary control run replaced `Workbench` with a nonsense string, `Zylorg`.

This shift was executed by dynamically injecting an alias mapping into the \texttt{LLMOrchestrator} mid-training, forcing the Commander to output weights for the alias rather than the canonical node name, which the orchestrator then reversed and fed to the PPO workers.

The initial hypothesis assumed the Commander would cleanly map "Furnace" to crafting success, while failing on "Zylorg". However, empirical results revealed a counter-intuitive phenomenon: \textbf{semantic over-confidence induced severe reward hacking, while semantic ignorance acted as a natural regularizer.}

\begin{figure}[H]
    \centering
    \includegraphics[width=0.48\textwidth]{../../plots/llm_dynamic_ablations_and_shifts/shift_subtask_Iron.png}
    \hfill
    \includegraphics[width=0.48\textwidth]{../../plots/llm_dynamic_ablations_and_shifts/shift_subtask_Sword.png}
    \caption{Subtask completion rates during semantic shifts. The `Furnace` shift traps workers in a local optimum, halting `Iron` and `Sword` progression, while the `Zylorg` shift organically recovers.}
    \label{fig:semantic_shift_subtasks}
\end{figure}

\begin{figure}[H]
    \centering
    \includegraphics[width=0.8\textwidth]{../../plots/llm_dynamic_ablations_and_shifts/shift_all_subtasks_bar.png}
    \caption{Maximum subtask completion rates for the Shift group. Furnace masters initial tasks (Wood, Stone, Pickaxe) but completely fails advanced tasks requiring exploration (Iron, Sword).}
    \label{fig:shift_all_subtasks}
\end{figure}

As seen in Figures~\ref{fig:semantic_shift_subtasks} and~\ref{fig:shift_all_subtasks}, the `Zylorg` control unexpectedly succeeded across almost all subtasks, achieving near-parity with the baseline's extrinsic reward. When the Commander encountered "Zylorg", it lacked a semantic prior. Consequently, it assigned it a balanced priority weight alongside other nodes, focusing its reasoning on recognizable end-goals like enemies.

\begin{hegquote}
\ttfamily
"reasoning": "Enemy defeat is the current bottleneck with no progress.",\\
"w\_wood": 0.8,\\
"w\_stone": 0.8,\\
"w\_zylorg": 0.9,\\
"w\_iron": 0.9,\\
"w\_enemy": 1.2
\end{hegquote}

Because \texttt{w\_zylorg} was not artificially inflated, the MAPPO workers, driven by the baseline extrinsic rewards for crafting a Sword, naturally explored the environment, discovered the `Zylorg` station was the required prerequisite, and solved the bottleneck via standard PPO exploration. The Commander's ignorance prevented it from interfering with the RL algorithm's foundational mechanics.

Conversely, the `Furnace` shift failed catastrophically on advanced tasks (`Iron`, `Sword`) due to reward hacking. The Commander correctly deduced zero-shot that a Furnace serves an identical semantic purpose to a crafting bench, as demonstrated in its reasoning logs:

\begin{hegquote}
\ttfamily
"reasoning": "Agents are not crafting Pickaxe or Sword, indicating a \\
bottleneck at the Workbench.",\\
"w\_furnace": 0.4,\\
"w\_iron": 0.2
\end{hegquote}

However, by mapping the Furnace to the crafting bottleneck, the Commander heavily skewed the Potential-Based Reward Shaping (PBRS) landscape. The MAPPO workers successfully crafted the Pickaxe (Figure~\ref{fig:shift_all_subtasks}), but subsequently discovered they could accumulate massive intrinsic returns simply by navigating back to the Furnace and idling there. They optimized for the dense shaping signal at the expense of the sparse extrinsic reward, halting all exploration for `Iron` and `Sword`. 

This finding highlights a critical vulnerability in neuro-symbolic MARL architectures: while LLMs possess exceptional zero-shot semantic mapping, coupling high-confidence oracle weights with dense gradient shaping can inadvertently trap low-level workers in local optima.
% ==========================================================================
% Chapter 8 — Conclusion
% ==========================================================================
\section{Conclusion: The Future of Neuro-Symbolic Control}
\label{ch:conclusion}

This thesis investigated the integration of Large Language Models (LLMs) with Multi-Agent Reinforcement Learning (MARL) to overcome exploration bottlenecks in sparse-reward environments. By structuring the environment as a strict Directed Acyclic Graph (DAG) of dependencies, we evaluated two distinct neuro-symbolic architectures: continuous Dynamic Reward Shaping (the "Mixing Board" approach) and discrete Hierarchical Reinforcement Learning (the "Meta-Controller" approach). 

The results demonstrate that while LLMs possess exceptional zero-shot semantic planning capabilities as high-level Commanders for low-level kinetic workers, integrating them introduces significant challenges regarding reward stationarity, semantic grounding, and temporal alignment.

\subsection{Review of Findings}
\label{sec:critical_assessment}

Our first major realization was the successful implementation of a dynamic Potential-Based Reward Shaping (PBRS) framework. We demonstrated that an LLM could actively monitor MARL agents and shift priority weights to guide exploration, successfully unblocking mid-game bottlenecks that a standard sparse-reward MAPPO baseline failed to solve. However, we proved the mathematical necessity of Exponential Moving Average (EMA) smoothing; without it, discrete jumps in the Commander's weight assignments induced spikes in Temporal Difference (TD) error variance, causing temporary regressions in the underlying PPO policies as well as being slower to solve the environment. 

The transition to a Hierarchical Reinforcement Learning (HRL) architecture represented a notable improvement in reliability. By allowing the Commander to dictate discrete macro-actions (Options) rather than adjusting the gravitational pull of rewards, the agents mastered the entire environment DAG, reaching 100\% completion on the final milestones. Through ablation, we established that QLoRA fine-tuning is highly beneficial for sample efficiency, enabling the agents to reach the final objective approximately $5.1\times$ faster than an un-finetuned base model.

Despite these successes, our experiments revealed vulnerabilities in neuro-symbolic MARL paradigms, highlighting a mismatch between an LLM's semantic priors and kinetic reality. We isolated two primary failure modes:

\begin{enumerate}
    \item We discovered that when an LLM possesses strong zero-shot knowledge of a concept (such as mapping a "Furnace" to crafting), it can confidently warp the shaping signal and trap the agents in a local optimum via reward hacking. Conversely, when presented with nonsense strings ("Zylorg"), its semantic ignorance acted as a natural regularizer, allowing the underlying PPO exploration mechanics to organically solve the bottleneck.
    \item Endowing the LLM with counterfactual reflection generated $2.3\times$ more staleness resets and delayed convergence. Because the LLM was fine-tuned on synthetic data where timeouts implied logical failures, it learned to doubt its own correct reasoning when faced with purely kinetic delays. This resulted in rapid, contradictory option switching that destabilized the PPO Critic, proving that reactive self-debugging is counterproductive if the model cannot distinguish between a logical error and physical friction.
\end{enumerate}

Finally, the computational overhead of running an LLM in the training loop remains a practical limitation for high-frequency control tasks, despite the use of quantized, localized models.

\subsection{Prospective Vision and Recommendations}
\label{sec:prospective_vision}

The findings of this thesis lay the groundwork for more robust neuro-symbolic architectures, carrying implications for complex coordination problems such as robotic swarm control, automated supply chain logistics, and advanced game AI. To realize this potential, future research must fundamentally bridge the gap between semantic logic and temporal delays. 

Rather than relying on simple text-based timeout triggers, Meta-Controllers should be equipped with spatio-temporal memory or multimodal state representations. This would allow the Commander to visually or temporally differentiate between a fundamentally flawed plan and a correct plan that merely requires the workers to spend more physical time navigating the grid.

Additionally, the stabilization of non-stationary objectives requires refinement. While our static KL-divergence guardrail ($\delta=0.015$) successfully stabilized the PPO workers during option transitions, implementing an \emph{adaptive trust region} could provide further robustness. Scaling the KL threshold dynamically based on the LLM's confidence or the specific Option type offers a more nuanced defense against objective shifts.

Perhaps the most promising avenue for future research is the concept of \emph{recursive self-improvement} through RL-driven self-alignment. Recent works like STaR \parencite{zelikman2022star} and Voyager \parencite{wang2023voyager} demonstrate that agentic LLMs can iteratively build skill libraries or generate reasoning traces that improve their own foundation. In our MARL context, the environment itself serves as an absolute verifier: if the LLM issues an Option and the PPO workers successfully achieve it, that trajectory constitutes ground-truth success. This successful reasoning trace can be distilled back into the LLM via DPO (Direct Preference Optimization), creating a closed-loop system where the LLM continuously refines its own kinetic understanding based on the physical successes and failures of its RL agents.

Finally, while the current architecture evaluates cooperative agents executing a shared Meta-Controller policy, a natural extension is the development of decentralized HRL. In this setting, multiple LLMs act as adversarial or dynamically allied faction leaders managing sub-groups of agents. Expanding into these complex diplomacy and resource-competition scenarios would thoroughly test the boundaries of MARL and AI alignment. 

This thesis demonstrates that while Large Language Models are profoundly powerful strategic directors for MARL agents, they must not be treated as infallible oracles. The path forward lies in designing architectures that safely balance the high-level semantic brilliance of language models with the uncompromising, low-level kinetic realities of reinforcement learning.

\newpage
\printbibliography[title={Bibliography}]


\newpage

\section*{Use of artificial intelligence-assisted tools}
\addcontentsline{toc}{section}{Use of artificial intelligence-assisted tools}

In the context of this work, the author states that he used artificial intelligence-assisted tools for the following purposes:

\begin{itemize}
    \item \textbf{Formal Improvements}: Assisted with spelling, syntax, phrasing, and report structuring throughout all chapters.\\
    \textit{AI tools used:} Gemini 3.1 Pro (via Antigravity Assistant)
    \item \textbf{Substantive Reflection}: Served as a collaborative sounding board for deep reflection during architectural design, debugging, and data interpretation across Chapters 4, 5, 6, and 7. The author remained the principal instigator and decision-maker for all methodologies and conclusions.\\
    \textit{AI tools used:} Gemini 3.1 Pro (via Antigravity Assistant)
    \item \textbf{Asset Generation}: Assisted in generating 2D pixel-art sprites and visual assets used in the custom MARL environment renderer (Chapter 4).\\
    \textit{AI tools used:} Gemini 1.5 Pro
\end{itemize}

\newpage
\newpage
\section*{Appendix 1: Core Algorithm Implementations}
\addcontentsline{toc}{section}{Appendix 1: Core Algorithm Implementations}

\begin{lstlisting}[style=code, language=Python, caption={DAG guardrail enforcement in \texttt{orchestrator.py}.}, label={lst:dag_guardrail}]
def apply_dag_guardrails(weights, metrics):
    guarded = dict(weights)
    # Iron requires Pickaxe
    if metrics.get('pickaxe', 0.0) < 1.0:
        guarded['w_iron'] = 0.0
    # Enemy and Gold require Bridge+Sword+Armor
    if (metrics.get('bridge', 0.0) < 1.0
        or metrics.get('sword', 0.0) < 1.0
        or metrics.get('armor', 0.0) < 1.0):
        guarded['w_enemy'] = 0.0
        guarded['w_gold'] = 0.0
    return guarded
\end{lstlisting}

\begin{lstlisting}[style=code, language=Python, caption={Vectorised PBRS computation across all environments in \texttt{orchestrator.py}.}, label={lst:pbrs_code}]
# Calculate L2 Norm distances for all agents simultaneously
dist_next = np.linalg.norm(pos_next - goal_targets, axis=2)
dist_prev = np.linalg.norm(pos_prev - goal_targets, axis=2)

MAX_DIST = 85.0
phi_next = MAX_DIST - dist_next
phi_prev = MAX_DIST - dist_prev

# Apply gamma discount, scaling, and the active mask
F = (gamma * phi_next - phi_prev) * R_LLM_SCALE * goal_active
\end{lstlisting}

\begin{lstlisting}[style=code, language=Python, caption={Per-agent intrinsic shaping in \texttt{hrl\_train\_loop.py}.}, label={lst:hrl_intrinsic}]
z0 = get_option_target_zone(a0_opt)
idle_mask_0 = (z0 == -1)
safe_z0 = np.where(idle_mask_0, 0, z0)

dist0_prev = np.linalg.norm(
    pos_prev[:, 0] - vec_env.zones[np.arange(n_envs), safe_z0],
    axis=1)
dist0_next = np.linalg.norm(
    vec_env.pos[:, 0] - vec_env.zones[np.arange(n_envs), safe_z0],
    axis=1)
phi0_prev = MAX_DIST - dist0_prev
phi0_next = MAX_DIST - dist0_next
shaping0 = np.where(
    idle_mask_0, 0.0,
    (0.99 * phi0_next) - phi0_prev
)
intrinsic_r[:, 0] = np.where(
    a0_success, 1.0, shaping0 * 0.05
)
\end{lstlisting}

\begin{lstlisting}[style=code, language=Python, caption={Composite reward in \texttt{hrl\_train\_loop.py}.}, label={lst:composite_reward}]
total_r = env_rewards + 0.5 * intrinsic_r
\end{lstlisting}

\begin{lstlisting}[style=code, language=Python, caption={Target KL guard in \texttt{hrl\_train\_loop.py}.}, label={lst:target_kl}]
with torch.no_grad():
    approx_kl = (-logratio).mean().item()
target_kl = 0.015
if approx_kl > target_kl:
    target_kl_reached = True
    break
\end{lstlisting}


\end{document}
